\ifnum\EditionMode=1\relax
\clearpage
\chapter*{Apparatus of corrections}
\addcontentsline{toc}{chapter}{Apparatus of corrections}
Each number identifies one particular occurrence in the corrected text. Fragments that cannot stand alone (an open delimiter or part of a diagram) are reproduced exactly in TeX notation. Transcription repairs shared by all three forms do not appear here.
\begingroup\small
\begin{description}
\item[\hypertarget{apparatus:EN-0001}{1}] Page \hyperlink{printpage:en:III}{III}\quad\texttt{SGA45-CORR-0001}
\par\noindent\textbf{Printed:} digression
\par\noindent\textbf{Corrected:} digressions
\par\noindent Plural series and certaines require digressions; print has singular.
\smallskip
\item[\hypertarget{apparatus:EN-0002}{2}] Page \hyperlink{printpage:en:10}{10}\quad\texttt{SGA45-CORR-0004}
\par\noindent\textbf{Printed:} flat is surjective
\par\noindent\textbf{Corrected:} flat and surjective
\par\noindent The printed est is a typographical substitution for conjunction et.
\smallskip
\item[\hypertarget{apparatus:EN-0003}{3}] Page \hyperlink{printpage:en:10}{10}\quad\texttt{S45-SOL-0027}
\par\noindent\textbf{Printed:} For every family
\par\noindent\textbf{Corrected:} For every finite family
\par\noindent The containment claim is valid for finite families; for infinite families it fails in an infinite free module.
\smallskip
\item[\hypertarget{apparatus:EN-0004}{4}] Page \hyperlink{printpage:en:12}{12}\quad\texttt{S45-SOL-0033}
\par\noindent\textbf{Printed:} \(M_U=\varphi^*M'\)
\par\noindent\textbf{Corrected:} \(M_V=\varphi^*M'\)
\par\noindent The pullback along V→U lives on V.
\smallskip
\item[\hypertarget{apparatus:EN-0005}{5}] Page \hyperlink{printpage:en:13}{13}\quad\texttt{S45-SOL-0035}
\par\noindent\textbf{Printed:} \(\partial_0\partial_1=\partial_0\partial_0\)
\par\noindent\textbf{Corrected:} \(\partial_1\partial_0=\partial_0\partial_0\)
\par\noindent The usual coface relation for right-to-left composition requires the reversed left indices.
\smallskip
\item[\hypertarget{apparatus:EN-0006}{6}] Page \hyperlink{printpage:en:14}{14}\quad\texttt{S45-SOL-0038}
\par\noindent\textbf{Printed:} \(v'\in V\)
\par\noindent\textbf{Corrected:} \(v'\in V'\)
\par\noindent The arbitrary basis vector is in V', not necessarily its invariant subspace V.
\smallskip
\item[\hypertarget{apparatus:EN-0007}{7}] Page \hyperlink{printpage:en:17}{17}\quad\texttt{SGA45-CORR-0006}
\par\noindent\textbf{Printed:} 6.5
\par\noindent\textbf{Corrected:} 6.6
\par\noindent Duplicated paragraph number; the next number is 6.6.
\smallskip
\item[\hypertarget{apparatus:EN-0008}{8}] Page \hyperlink{printpage:en:20}{20}\quad\texttt{S45-SOL-0057}
\par\noindent\textbf{Printed:} \hyperref[I:1-6-6]{I.6.5}
\par\noindent\textbf{Corrected:} \hyperref[I:1-6-6]{I.6.6}
\par\noindent Reference follows the corrected torsor paragraph number.
\smallskip
\item[\hypertarget{apparatus:EN-0009}{9}] Page \hyperlink{printpage:en:24}{24}\quad\texttt{S45-SOL-0085}
\par\noindent\textbf{Printed:} the spectrum of the strict localization of $Y$ at $y$
\par\noindent\textbf{Corrected:} the spectrum of the strict localization of $Y$ at $\bar y$
\par\noindent Strict localization is at the chosen geometric point.
\smallskip
\item[\hypertarget{apparatus:EN-0010}{10}] Page \hyperlink{printpage:en:28}{28}\quad\texttt{S45-SOL-0094}
\par\noindent\textbf{Printed:} of $K$ for which $A_K$
\par\noindent\textbf{Corrected:} of $K$ for which $A_{K'}$
\par\noindent The indicated extension of scalars is to K′.
\smallskip
\item[\hypertarget{apparatus:EN-0011}{11}] Page \hyperlink{printpage:en:30}{30}\quad\texttt{S45-SOL-0103}
\par\noindent\textbf{Printed:} \(\h^1\left(G,(\bar K\otimes_K K')^*/\bar k^*\right)\)
\par\noindent\textbf{Corrected:} \(\h^1\left(G,(\bar K\otimes_K K')^*/\bar K^*\right)\)
\par\noindent The closure is that of K; lowercase k is undefined.
\smallskip
\item[\hypertarget{apparatus:EN-0012}{12}] Page \hyperlink{printpage:en:32}{32}\quad\texttt{S45-SOL-0114}
\par\noindent\textbf{Printed:} \(g(X_{ij})\)
\par\noindent\textbf{Corrected:} \(g(U_{ij})\)
\par\noindent The variables introduced in the previous sentence are U\_\{ij\}.
\smallskip
\item[\hypertarget{apparatus:EN-0013}{13}] Page \hyperlink{printpage:en:32}{32}\quad\texttt{SGA45-CORR-0011B}
\par\noindent\textbf{Printed:} \(N_{K/k}\)
\par\noindent\textbf{Corrected:} \(N_{K/k(X)}\)
\par\noindent Only K/k(X) is the indicated finite extension.
\smallskip
\item[\hypertarget{apparatus:EN-0014}{14}] Page \hyperlink{printpage:en:45}{45}\quad\texttt{S45-SOL-0146}
\par\noindent\textbf{Printed:} form $\prod_{p_{i*}}C_i$
\par\noindent\textbf{Corrected:} form $\prod p_{i*}C_i$
\par\noindent The product is of the direct images p\_\{i*\}C\_i, as in (3.2)(ii).
\smallskip
\item[\hypertarget{apparatus:EN-0015}{15}] Page \hyperlink{printpage:en:45}{45}\quad\texttt{S45-ASTRA-A-0021}
\par\noindent\textbf{Printed:} is an integer $\geqslant0$
\par\noindent\textbf{Corrected:} is an integer $>0$
\par\noindent The reduction uses torsion coefficients Z/n and then reduces n to a prime power; n=0 would give Z and is outside that argument.
\smallskip
\item[\hypertarget{apparatus:EN-0016}{16}] Page \hyperlink{printpage:en:49}{49}\quad\texttt{SGA45-CORR-0015}
\par\noindent\textbf{Printed:} the 
\par\noindent\textbf{Corrected:} \emph{\textnormal{[omission]}}
\par\noindent The preceding page already ends with le; remove only the duplicated article in corrected/reading.
\smallskip
\item[\hypertarget{apparatus:EN-0017}{17}] Page \hyperlink{printpage:en:49}{49}\quad\texttt{S45-SOL-0154}
\par\noindent\textbf{Printed:} \(\eR^q f_! =\)
\par\noindent\textbf{Corrected:} \(\eR^q f_! \sF =\)
\par\noindent The displayed equality defines the value at F.
\smallskip
\item[\hypertarget{apparatus:EN-0018}{18}] Page \hyperlink{printpage:en:49}{49}\quad\texttt{SGA45-CORR-0014}
\par\noindent\textbf{Printed:} \(\eR^q f_*(j_! \sF)\)
\par\noindent\textbf{Corrected:} \(\eR^q \bar f_*(j_! \sF)\)
\par\noindent The argument is a sheaf on the compactification, whose map is barred f.
\smallskip
\item[\hypertarget{apparatus:EN-0019}{19}] Page \hyperlink{printpage:en:54}{54}\quad\texttt{SGA45-CORR-0019}
\par\noindent\textbf{Printed:} \(\h^0\left(X_t^x,\dZ/n\right)\)
\par\noindent\textbf{Corrected:} \(\h^0\left(\widetilde X_t^x,\dZ/n\right)\)
\par\noindent The expression concerns the vanishing-cycle variety already defined with a tilde.
\smallskip
\item[\hypertarget{apparatus:EN-0020}{20}] Page \hyperlink{printpage:en:56}{56}\quad\texttt{S45-ASTRA-A-0020}
\par\noindent\textbf{Printed:} \(\h^*(X,\varepsilon_*' \dZ/n)\)
\par\noindent\textbf{Corrected:} \(\h^*(\widetilde X,\varepsilon_*' \dZ/n)\)
\par\noindent The direct image by epsilon-prime is a sheaf on X-tilde in the preceding Cartesian diagram.
\smallskip
\item[\hypertarget{apparatus:EN-0021}{21}] Page \hyperlink{printpage:en:61}{61}\quad\texttt{S45-SOL-0209}
\par\noindent\textbf{Printed:} \(\widehat X=A\llbracket T\rrbracket\)
\par\noindent\textbf{Corrected:} \(\widehat X=\spec A\llbracket T\rrbracket\)
\par\noindent The completed base is a scheme; Spec is missing from the printed equality.
\smallskip
\item[\hypertarget{apparatus:EN-0022}{22}] Page \hyperlink{printpage:en:64}{64}\quad\texttt{S45-SOL-0217}
\par\noindent\textbf{Printed:} \(\dA_S^{n-1}\)
\par\noindent\textbf{Corrected:} \(\dA_S^{N-1}\)
\par\noindent N is the relative dimension; n is the torsion modulus.
\smallskip
\item[\hypertarget{apparatus:EN-0023}{23}] Page \hyperlink{printpage:en:74}{74}\quad\texttt{S45-ASTRA-A-0029}
\par\noindent\textbf{Printed:} 1959
\par\noindent\textbf{Corrected:} 1969
\par\noindent The original journal record gives volume36 (1969), DOI10.1007/BF02684596; the printed bibliography has the wrong decade. Verified at https://www.numdam.org/item/PMIHES\_1969\_\_36\_\_23\_0/.
\smallskip
\item[\hypertarget{apparatus:EN-0024}{24}] Page \hyperlink{printpage:en:74}{74}\quad\texttt{S45-ASTRA-A-0030}
\par\noindent\textbf{Printed:} 25--58
\par\noindent\textbf{Corrected:} 23--58
\par\noindent The journal record and the article first page give23–58, not25–58. DOI10.1007/BF02684596; https://www.numdam.org/item/PMIHES\_1969\_\_36\_\_23\_0/.
\smallskip
\item[\hypertarget{apparatus:EN-0025}{25}] Page \hyperlink{printpage:en:74}{74}\quad\texttt{S45-ASTRA-A-0033}
\par\noindent\textbf{Printed:} 266
\par\noindent\textbf{Corrected:} 256
\par\noindent The same work is correctly numbered256 in Ariane p3; its primary Bourbaki record confirms256: https://www.numdam.org/item/SB\_1962-1964\_\_8\_\_189\_0/.
\smallskip
\item[\hypertarget{apparatus:EN-0026}{26}] Page \hyperlink{printpage:en:75}{75}\quad\texttt{S45-ASTRA-A-0031}
\par\noindent\textbf{Printed:} 1--12
\par\noindent\textbf{Corrected:} 1--42
\par\noindent The original journal record gives Annales de l’Institut Fourier6 (1956),1–42, DOI10.5802/aif.59; https://www.numdam.org/item/?id=AIF\_1956\_\_6\_\_1\_0. The printed final page12 is wrong.
\smallskip
\item[\hypertarget{apparatus:EN-0027}{27}] Page \hyperlink{printpage:en:85}{85}\quad\texttt{S45-SOL-0309}
\par\noindent\textbf{Printed:} $\dZ_\ell$-sheaf on $Y$,
\par\noindent\textbf{Corrected:} $\dZ_\ell$-sheaf on $X$,
\par\noindent For the declared f:X→Y, the input to R\textasciicircum{}q f\_! must be a sheaf on X; the print visibly has Y in that hypothesis. The output remains on Y.
\smallskip
\item[\hypertarget{apparatus:EN-0028}{28}] Page \hyperlink{printpage:en:92}{92}\quad\texttt{S45-SOL-0342}
\par\noindent\textbf{Printed:} $\gr_F^p(\Lambda)$ belong
\par\noindent\textbf{Corrected:} $\gr_F^p(K)$ belong
\par\noindent The source visibly prints Lambda in the graded object. The surrounding pair (K,F), preceding graded terms, and essential-image criterion require the graded pieces of K.
\smallskip
\item[\hypertarget{apparatus:EN-0029}{29}] Page \hyperlink{printpage:en:93}{93}\quad\texttt{S45-SOL-0345}
\par\noindent\textbf{Printed:} \(\tor_k(N,K^r/\im d)\)
\par\noindent\textbf{Corrected:} \(\tor_k(N,K^{-r}/\im d)\)
\par\noindent The visible exponent r lacks a minus sign; the preceding free resolution starts in degree -r and computes Tor with its cokernel.
\smallskip
\item[\hypertarget{apparatus:EN-0030}{30}] Page \hyperlink{printpage:en:93}{93}\quad\texttt{S45-SOL-0346}
\par\noindent\textbf{Printed:} the module $K^r/\im d$
\par\noindent\textbf{Corrected:} the module $K^{-r}/\im d$
\par\noindent This is the cokernel of the resolution in degree -r; the scan repeats the missing minus sign.
\smallskip
\item[\hypertarget{apparatus:EN-0031}{31}] Page \hyperlink{printpage:en:93}{93}\quad\texttt{S45-SOL-0347}
\par\noindent\textbf{Printed:} \(0 \to K^r/\im(d) \to K^{r+1} \to K^{r+2}\)
\par\noindent\textbf{Corrected:} \(0 \to K^{-r}/\im(d) \to K^{-r+1} \to K^{-r+2}\)
\par\noindent All three displayed exponents visibly use positive r; the truncation must start at -r, consistently with the resolution and Tor-dimension bound.
\smallskip
\item[\hypertarget{apparatus:EN-0032}{32}] Page \hyperlink{printpage:en:94}{94}\quad\texttt{S45-SOL-0351}
\par\noindent\textbf{Printed:} If a complex of sheaves
\par\noindent\textbf{Corrected:} If a complex $\sK$ of sheaves
\par\noindent The print omits the name of the complex but immediately uses K in H\textasciicircum{}i(K) and K′→K; supply the missing name in corrected forms.
\smallskip
\item[\hypertarget{apparatus:EN-0033}{33}] Page \hyperlink{printpage:en:94}{94}\quad\texttt{S45-SOL-0353}
\par\noindent\textbf{Printed:} construct $\sK^n$ flat and constructible and $v:\sK^n\to \sA$
\par\noindent\textbf{Corrected:} construct ${\sK'}^n$ flat and constructible and $v:{\sK'}^n\to \sA$
\par\noindent The scan twice gives unprimed K\textasciicircum{}n, although the induction constructs the terms of K′.
\smallskip
\item[\hypertarget{apparatus:EN-0034}{34}] Page \hyperlink{printpage:en:95}{95}\quad\texttt{S45-SOL-0355}
\par\noindent\textbf{Printed:} Let $F:X\to Y$ be a separated morphism
\par\noindent\textbf{Corrected:} Let $f:X\to Y$ be a separated morphism
\par\noindent Capital F in the theorem is inconsistent with the subsequent functor f; the diagram itself correctly has lowercase barred f and is not changed.
\smallskip
\item[\hypertarget{apparatus:EN-0035}{35}] Page \hyperlink{printpage:en:95}{95}\quad\texttt{S45-SOL-0357}
\par\noindent\textbf{Printed:} $\eR f_!\sK\in \ob \eD^b(X,\Lambda)$
\par\noindent\textbf{Corrected:} $\eR f_!\sK\in \ob \eD^b(Y,\Lambda)$
\par\noindent The scan prints X, but pushforward under f:X→Y is a complex on Y, as the theorem states.
\smallskip
\item[\hypertarget{apparatus:EN-0036}{36}] Page \hyperlink{printpage:en:95}{95}\quad\texttt{S45-SOL-0359}
\par\noindent\textbf{Printed:} $\eR f_! = \eR f_*\circ j_!$
\par\noindent\textbf{Corrected:} $\eR f_! = \eR \bar f_*\circ j_!$
\par\noindent The compactification uses proper map bar f. The bottom line drops the bar, unlike the definition above; corrected forms restore it.
\smallskip
\item[\hypertarget{apparatus:EN-0037}{37}] Page \hyperlink{printpage:en:96}{96}\quad\texttt{S45-SOL-0362}
\par\noindent\textbf{Printed:} when $\Lambda$ is free of finite type
\par\noindent\textbf{Corrected:} when $L$ is free of finite type
\par\noindent The free right module being varied is L; Lambda is the coefficient ring. The following passage-to-the-limit argument concerns L.
\smallskip
\item[\hypertarget{apparatus:EN-0038}{38}] Page \hyperlink{printpage:en:97}{97}\quad\texttt{S45-SOL-0370}
\par\noindent\textbf{Printed:} \(\sK\in\ob_{\text{ctf}}(X,\Lambda)\)
\par\noindent\textbf{Corrected:} \(\sK\in\ob\eD_\text{ctf}(X,\Lambda)\)
\par\noindent The scan omits the category letter D after Ob; the object must belong to the derived category with constructible finite-Tor-dimension cohomology.
\smallskip
\item[\hypertarget{apparatus:EN-0039}{39}] Page \hyperlink{printpage:en:98}{98}\quad\texttt{S45-SOL-0379}
\par\noindent\textbf{Printed:} $K_n=\eR\Gamma_c(X,\sF\otimes\dZ/\ell^n)$
\par\noindent\textbf{Corrected:} $K_n=\eR\Gamma_c(X,\sF\otimes\dZ/\ell^{n+1})$
\par\noindent The coefficient exponent n contradicts the following category D\_parf(Z/l\textasciicircum{}\{n+1\}), formula4.11.1, and trace computation on p.99; restore n+1 only in corrected forms.
\smallskip
\item[\hypertarget{apparatus:EN-0040}{40}] Page \hyperlink{printpage:en:99}{99}\quad\texttt{S45-ASTRA-R-0006}
\par\noindent\textbf{Printed:} with the sign convention of .1
\par\noindent\textbf{Corrected:} with the sign convention of \ref{II:eq:4-3-1}
\par\noindent The incomplete printed reference .1 refers to the alternating trace definition(4.3.1); Codex corrected restores the same reference.
\smallskip
\item[\hypertarget{apparatus:EN-0041}{41}] Page \hyperlink{printpage:en:100}{100}\quad\texttt{S45-SOL-0386}
\par\noindent\textbf{Printed:} formula (4.10.1) holds
\par\noindent\textbf{Corrected:} formula \eqref{II:eq:4-10} holds
\par\noindent The trace display is visibly numbered(4.10) on p.96; the printed cross-reference(4.10.1) has a spurious final component.
\smallskip
\item[\hypertarget{apparatus:EN-0042}{42}] Page \hyperlink{printpage:en:105}{105}\quad\texttt{S45-SOL-0419}
\par\noindent\textbf{Printed:} The complex of $\Lambda[H]$-modules
\par\noindent\textbf{Corrected:} The complex of $A[H]$-modules
\par\noindent RΓ\_c(U,f\_*A) has coefficient ring A[H], as defined on p.104; Lambda appears only after tensoring with M.
\smallskip
\item[\hypertarget{apparatus:EN-0043}{43}] Page \hyperlink{printpage:en:105}{105}\quad\texttt{S45-SOL-0420}
\par\noindent\textbf{Printed:} a complex of $\Lambda[G^+]$-modules
\par\noindent\textbf{Corrected:} a complex of $A[G^+]$-modules
\par\noindent The same A[H]-complex with Frobenius action is an A[G+]-complex; the printed Lambda is inconsistent with its coefficients.
\smallskip
\item[\hypertarget{apparatus:EN-0044}{44}] Page \hyperlink{printpage:en:105}{105}\quad\texttt{S45-SOL-0421}
\par\noindent\textbf{Printed:} Formulas \ref{II:5-10}
\par\noindent\textbf{Corrected:} Formulas \ref{II:5-11}
\par\noindent The coinvariant trace formula used is proposition5.11; paragraph5.10 introduces the notation but contains no trace identity.
\smallskip
\item[\hypertarget{apparatus:EN-0045}{45}] Page \hyperlink{printpage:en:110}{110}\quad\texttt{S45-ASTRA-R-0023}
\par\noindent\textbf{Printed:} Artin-Scheier
\par\noindent\textbf{Corrected:} Artin--Schreier
\par\noindent Printed eponym misspelled.
\smallskip
\item[\hypertarget{apparatus:EN-0046}{46}] Page \hyperlink{printpage:en:112}{112}\quad\texttt{S45-SOL-0449}
\par\noindent\textbf{Printed:} \(\det(1-u t,\gr_F^p(F))\)
\par\noindent\textbf{Corrected:} \(\det(1-u t,\gr_F^p(K))\)
\par\noindent The print has F as the argument; the filtered complex is K, as stated immediately above.
\smallskip
\item[\hypertarget{apparatus:EN-0047}{47}] Page \hyperlink{printpage:en:113}{113}\quad\texttt{S45-SOL-0450}
\par\noindent\textbf{Printed:} \(s(\Gamma^i K'\otimes \Gamma^i K'')\)
\par\noindent\textbf{Corrected:} \(s(\Gamma^i K'\otimes \Gamma^j K'')\)
\par\noindent The second tensor exponent must be j in the i+j=n summand.
\smallskip
\item[\hypertarget{apparatus:EN-0048}{48}] Page \hyperlink{printpage:en:113}{113}\quad\texttt{S45-SOL-0452}
\par\noindent\textbf{Printed:} $\Gamma^i K'\otimes \Gamma^i K''$, and
\par\noindent\textbf{Corrected:} $\Gamma^i K'\otimes \Gamma^j K''$, and
\par\noindent The double-complex second tensor exponent must likewise be j.
\smallskip
\item[\hypertarget{apparatus:EN-0049}{49}] Page \hyperlink{printpage:en:114}{114}\quad\texttt{S45-SOL-0458}
\par\noindent\textbf{Printed:} \(=_{\text{(\ref{III:1-10})}}\)
\par\noindent\textbf{Corrected:} \(=_{\text{(\ref{III:1-11})}}\)
\par\noindent The first equality uses the filtration lemma 1.10; the second uses the shifted one-term computation of 1.11.
\smallskip
\item[\hypertarget{apparatus:EN-0050}{50}] Page \hyperlink{printpage:en:118}{118}\quad\texttt{S45-SOL-0480}
\par\noindent\textbf{Printed:} \(\sum_n (\tr\left(F^*,\Gamma^n \Gamma(X,\sK)\right) T^n\)
\par\noindent\textbf{Corrected:} \(\sum_n \tr\left(F^*,\Gamma^n \Gamma(X,\sK)\right) T^n\)
\par\noindent An unmatched opening parenthesis occurs immediately after the printed summation sign.
\smallskip
\item[\hypertarget{apparatus:EN-0051}{51}] Page \hyperlink{printpage:en:118}{118}\quad\texttt{S45-SOL-0481}
\par\noindent\textbf{Printed:} \(\tr\left(F^*,\Gamma(\operatorname{Sym}^n(X),\Gamma_\text{ext}^n(\sK)\right)\)
\par\noindent\textbf{Corrected:} \(\tr\left(F^*,\Gamma(\operatorname{Sym}^n(X),\Gamma_\text{ext}^n(\sK))\right)\)
\par\noindent The printed formula is one closing parenthesis short; the section functor Gamma needs closing before the trace.
\smallskip
\item[\hypertarget{apparatus:EN-0052}{52}] Page \hyperlink{printpage:en:119}{119}\quad\texttt{S45-SOL-0483}
\par\noindent\textbf{Printed:} \(\tr\left(F^*,\eR\Gamma_c(\operatorname{Sym}^n(X,\Gamma_\text{ext}^n \sK)\right)\)
\par\noindent\textbf{Corrected:} \(\tr\left(F^*,\eR\Gamma_c(\operatorname{Sym}^n(X),\Gamma_\text{ext}^n \sK)\right)\)
\par\noindent The print omits the parenthesis closing Sym\textasciicircum{}n(X).
\smallskip
\item[\hypertarget{apparatus:EN-0053}{53}] Page \hyperlink{printpage:en:120}{120}\quad\texttt{S45-ASTRA-R-0012}
\par\noindent\textbf{Printed:} Theory of Artin-Scheier
\par\noindent\textbf{Corrected:} Theory of Artin--Schreier
\par\noindent Printed eponym misspelled in section heading.
\smallskip
\item[\hypertarget{apparatus:EN-0054}{54}] Page \hyperlink{printpage:en:120}{120}\quad\texttt{S45-ASTRA-R-0013}
\par\noindent\textbf{Printed:} Artin-Scheier theory \PrintCite{SGA~4 IX 3.5}{\cite[IX 3.5]{sga4}}
\par\noindent\textbf{Corrected:} Artin--Schreier theory \PrintCite{SGA~4 IX 3.5}{\cite[IX 3.5]{sga4}}
\par\noindent Printed eponym misspelled in text.
\smallskip
\item[\hypertarget{apparatus:EN-0055}{55}] Page \hyperlink{printpage:en:120}{120}\quad\texttt{S45-SOL-0490}
\par\noindent\textbf{Printed:} \(\Phi(I^n \sG')\subset I^{p^n}\Phi(\sG')\)
\par\noindent\textbf{Corrected:} \(\Phi(I^n \sG')\subset I^{p n}\Phi(\sG')\)
\par\noindent p-semilinearity multiplies the ideal exponent by p; the printed raised n is a typo.
\smallskip
\item[\hypertarget{apparatus:EN-0056}{56}] Page \hyperlink{printpage:en:123}{123}\quad\texttt{S45-ASTRA-R-0015}
\par\noindent\textbf{Printed:} Artin-Scheier theory provides
\par\noindent\textbf{Corrected:} Artin--Schreier theory provides
\par\noindent Printed eponym misspelled.
\smallskip
\item[\hypertarget{apparatus:EN-0057}{57}] Page \hyperlink{printpage:en:123}{123}\quad\texttt{S45-SOL-0501}
\par\noindent\textbf{Printed:} (cf. \ref{III:3-6}b)
\par\noindent\textbf{Corrected:} (cf. \ref{III:3-3}c)
\par\noindent Section3.6 has no part b. The surjectivity argument is explicitly proved in3.3c.
\smallskip
\item[\hypertarget{apparatus:EN-0058}{58}] Page \hyperlink{printpage:en:124}{124}\quad\texttt{S45-SOL-0506}
\par\noindent\textbf{Printed:} \(\ker(\Phi-1)\otimes_{\dF_q}\dF \iso\)
\par\noindent\textbf{Corrected:} \(\ker(\Phi-1)\otimes_{\dF_p}\dF \iso\)
\par\noindent The Phi-fixed vectors form an F\_p-space. The next printed trace chain correctly uses F\_p.
\smallskip
\item[\hypertarget{apparatus:EN-0059}{59}] Page \hyperlink{printpage:en:125}{125}\quad\texttt{S45-SOL-0509}
\par\noindent\textbf{Printed:} each side of (3.1.1)
\par\noindent\textbf{Corrected:} each side of \eqref{III:eq:4-1-1}
\par\noindent The cited trace formula is4.1.1, not3.1.1.
\smallskip
\item[\hypertarget{apparatus:EN-0060}{60}] Page \hyperlink{printpage:en:125}{125}\quad\texttt{S45-SOL-0512}
\par\noindent\textbf{Printed:} the two sides of (3.1.1)
\par\noindent\textbf{Corrected:} the two sides of \eqref{III:eq:4-1-1}
\par\noindent The two trace expressions belong to4.1.1.
\smallskip
\item[\hypertarget{apparatus:EN-0061}{61}] Page \hyperlink{printpage:en:125}{125}\quad\texttt{S45-SOL-0513}
\par\noindent\textbf{Printed:} If (3.1.1) holds
\par\noindent\textbf{Corrected:} If \eqref{III:eq:4-1-1} holds
\par\noindent The reduction assumes trace formula4.1.1.
\smallskip
\item[\hypertarget{apparatus:EN-0062}{62}] Page \hyperlink{printpage:en:126}{126}\quad\texttt{S45-ASTRA-R-0018}
\par\noindent\textbf{Printed:} Woodshole fixed point formula
\par\noindent\textbf{Corrected:} Woods Hole fixed point formula
\par\noindent Correct the printed spelling of the Woods Hole name.
\smallskip
\item[\hypertarget{apparatus:EN-0063}{63}] Page \hyperlink{printpage:en:127}{127}\quad\texttt{S45-ASTRA-R-0019}
\par\noindent\textbf{Printed:} \(u^* F^*\)
\par\noindent\textbf{Corrected:} \(u F^*\)
\par\noindent The trace is of the composition defined immediately before4.4.3; no extra star on u.
\smallskip
\item[\hypertarget{apparatus:EN-0064}{64}] Page \hyperlink{printpage:en:127}{127}\quad\texttt{S45-SOL-0520}
\par\noindent\textbf{Printed:} \(\frac{\tr(u_x,\sG_x)}\)
\par\noindent\textbf{Corrected:} \(\frac{\tr(u_x,\bar\sG_x)}\)
\par\noindent The trace must use the extended sheaf on all of barS.
\smallskip
\item[\hypertarget{apparatus:EN-0065}{65}] Page \hyperlink{printpage:en:127}{127}\quad\texttt{S45-ASTRA-R-0020}
\par\noindent\textbf{Printed:} the endomorphism $u_x$ of $\sG_x$
\par\noindent\textbf{Corrected:} the endomorphism $u_x$ of $\bar\sG_x$
\par\noindent The extended sheaf is required on the boundary; the scan omits its bar at this occurrence.
\smallskip
\item[\hypertarget{apparatus:EN-0066}{66}] Page \hyperlink{printpage:en:137}{137}\quad\texttt{S45-SOL-0552}
\par\noindent\textbf{Printed:} \(V^\vee\otimes V\otimes W^\vee\otimes W\to k\)
\par\noindent\textbf{Corrected:} \(V^\vee\otimes V\otimes W^\vee\otimes W\to A\)
\par\noindent The ground field in the surrounding definition is A, including the evaluation maps1.3.1 and1.3.4; k has not been introduced.
\smallskip
\item[\hypertarget{apparatus:EN-0067}{67}] Page \hyperlink{printpage:en:141}{141}\quad\texttt{S45-SOL-0572}
\par\noindent\textbf{Printed:} \(\eR k^! \dZ/n=\eR k_1^! \eR i^!\)
\par\noindent\textbf{Corrected:} \(\eR k^! \dZ/n=\eR k_1^! \eR i^!\dZ/n\)
\par\noindent The right side lacks the coefficient object, while the left is the functor applied to Z/n.
\smallskip
\item[\hypertarget{apparatus:EN-0068}{68}] Page \hyperlink{printpage:en:142}{142}\quad\texttt{S45-SOL-0577}
\par\noindent\textbf{Printed:} \(\eR^{2 c+q} i^!(g^*\sF)=0\Leftrightarrow \eR^{-2 d+q}f^!(\sF)\)
\par\noindent\textbf{Corrected:} \(\eR^{2 c+q} i^!(g^*\sF)=0\Leftrightarrow \eR^{-2 d+q}f^!(\sF)=0\)
\par\noindent The right side of the vanishing equivalence requires =0, omitted in print.
\smallskip
\item[\hypertarget{apparatus:EN-0069}{69}] Page \hyperlink{printpage:en:144}{144}\quad\texttt{S45-SOL-0585}
\par\noindent\textbf{Printed:} \(\h^0\left(Y,\eR^{-2 d}f^!\dZ/n\right)\)
\par\noindent\textbf{Corrected:} \(\h^0\left(Y,\eR^{-2 d}f^!\dZ/n(-d)\right)\)
\par\noindent The preceding Hom term has twist -d, which survives into the final degree-zero cohomology group.
\smallskip
\item[\hypertarget{apparatus:EN-0070}{70}] Page \hyperlink{printpage:en:144}{144}\quad\texttt{S45-SOL-0588}
\par\noindent\textbf{Printed:} \(\eR g^! \dZ/n=\dZ/n(-N)[-2 N]\)
\par\noindent\textbf{Corrected:} \(\eR g^! \dZ/n=\dZ/n(N)[2 N]\)
\par\noindent Smooth purity for relative dimension N gives positive twist and shift, as also used on p142.
\smallskip
\item[\hypertarget{apparatus:EN-0071}{71}] Page \hyperlink{printpage:en:145}{145}\quad\texttt{S45-ASTRA-R-0031}
\par\noindent\textbf{Printed:} \(\otimes \eR^{2 d} g_!\dZ/n(d)\)
\par\noindent\textbf{Corrected:} \(\otimes \eR^{2 d} f_!\dZ/n(d)\)
\par\noindent The input pushforward must come from Y via f to match the trace relation.
\smallskip
\item[\hypertarget{apparatus:EN-0072}{72}] Page \hyperlink{printpage:en:145}{145}\quad\texttt{S45-SOL-0593}
\par\noindent\textbf{Printed:} where the cup product can be computed in $\h^*(X,\dZ/n)$.
\par\noindent\textbf{Corrected:} where the cup product can be computed in $\h^*(X,\dZ/n(c))$.
\par\noindent The cycle class has twist c; cup with untwisted v retains this twist.
\smallskip
\item[\hypertarget{apparatus:EN-0073}{73}] Page \hyperlink{printpage:en:146}{146}\quad\texttt{S45-ASTRA-R-0033}
\par\noindent\textbf{Printed:} \PrintCite{SGA~4 XVII \S2}{\cite[XVII \S2]{sga4}}
\par\noindent\textbf{Corrected:} \PrintCite{SGA~4 XVIII \S2}{\cite[XVIII \S2]{sga4}}
\par\noindent The printed XVII differs from the systematic trace construction of XVIII2 referred to in the continuation.
\smallskip
\item[\hypertarget{apparatus:EN-0074}{74}] Page \hyperlink{printpage:en:146}{146}\quad\texttt{S45-SOL-0600}
\par\noindent\textbf{Printed:} \(u:Z\to \dA_S^d\)
\par\noindent\textbf{Corrected:} \(u:X\to \dA_S^d\)
\par\noindent The composite ui requires u to have source X because i:Y→X.
\smallskip
\item[\hypertarget{apparatus:EN-0075}{75}] Page \hyperlink{printpage:en:147}{147}\quad\texttt{S45-SOL-0603}
\par\noindent\textbf{Printed:} \(\h_X^{2 c}\left(Z,\dZ/n(c)\right)\)
\par\noindent\textbf{Corrected:} \(\h_Y^{2 c}\left(X,\dZ/n(c)\right)\)
\par\noindent The lemma concerns the class of Y supported in X; Z is not an ambient scheme here.
\smallskip
\item[\hypertarget{apparatus:EN-0076}{76}] Page \hyperlink{printpage:en:147}{147}\quad\texttt{S45-SOL-0605}
\par\noindent\textbf{Printed:} \(\hom_{S'}\left(\eR^{2 d'} f_!' \dZ/n,\dZ/n\right)\)
\par\noindent\textbf{Corrected:} \(\hom_{S'}\left(\eR^{2 d'} f_!' \dZ/n(d'),\dZ/n\right)\)
\par\noindent The trace over S requires the d source twist.
\smallskip
\item[\hypertarget{apparatus:EN-0077}{77}] Page \hyperlink{printpage:en:147}{147}\quad\texttt{S45-SOL-0606}
\par\noindent\textbf{Printed:} \(\eR g_!'\dZ/n(N)[2 N]\)
\par\noindent\textbf{Corrected:} \(\eR g_!\dZ/n(N)[2 N]\)
\par\noindent The bottom row pushes forward from X to S via g; the prime is erroneous.
\smallskip
\item[\hypertarget{apparatus:EN-0078}{78}] Page \hyperlink{printpage:en:147}{147}\quad\texttt{S45-SOL-0605}
\par\noindent\textbf{Printed:} \(\hom_S\left(\eR^{2 d} f_!\dZ/n,\dZ/n\right)\)
\par\noindent\textbf{Corrected:} \(\hom_S\left(\eR^{2 d} f_!\dZ/n(d),\dZ/n\right)\)
\par\noindent The trace over S requires the d source twist.
\smallskip
\item[\hypertarget{apparatus:EN-0079}{79}] Page \hyperlink{printpage:en:149}{149}\quad\texttt{S45-SOL-0612}
\par\noindent\textbf{Printed:} \(\hom_S\left(\eR^{2 d'} f_!' \dZ/n,\dZ/n\right)\)
\par\noindent\textbf{Corrected:} \(\hom_S\left(\eR^{2 d'} f_!' \dZ/n(d'),\dZ/n\right)\)
\par\noindent The product trace requires its source twist d.
\smallskip
\item[\hypertarget{apparatus:EN-0080}{80}] Page \hyperlink{printpage:en:149}{149}\quad\texttt{S45-SOL-0612}
\par\noindent\textbf{Printed:} \(\hom_S\left(\eR^{2 d}f_!\dZ/n,\dZ/n\right)\)
\par\noindent\textbf{Corrected:} \(\hom_S\left(\eR^{2 d}f_!\dZ/n(d),\dZ/n\right)\)
\par\noindent The product trace requires its source twist d.
\smallskip
\item[\hypertarget{apparatus:EN-0081}{81}] Page \hyperlink{printpage:en:149}{149}\quad\texttt{S45-SOL-0613}
\par\noindent\textbf{Printed:} with $X'=X'$
\par\noindent\textbf{Corrected:} with $Y'=X'$
\par\noindent The cycle in the first factor must be all of X-prime; the printed tautology does not supply that cycle.
\smallskip
\item[\hypertarget{apparatus:EN-0082}{82}] Page \hyperlink{printpage:en:149}{149}\quad\texttt{S45-SOL-0615}
\par\noindent\textbf{Printed:} \(\cl(Y')=1\in \h^0(X,\dZ/n)\)
\par\noindent\textbf{Corrected:} \(\cl(Y')=1\in \h^0(X',\dZ/n)\)
\par\noindent The unit class belongs to the first factor X-prime.
\smallskip
\item[\hypertarget{apparatus:EN-0083}{83}] Page \hyperlink{printpage:en:150}{150}\quad\texttt{S45-SOL-0619}
\par\noindent\textbf{Printed:} By (iv), one therefore has
\par\noindent\textbf{Corrected:} By (iii), one therefore has
\par\noindent Theorem2.3.8 has only three assertions; its product assertion is(iii).
\smallskip
\item[\hypertarget{apparatus:EN-0084}{84}] Page \hyperlink{printpage:en:151}{151}\quad\texttt{S45-SOL-0621}
\par\noindent\textbf{Printed:} $2(q-N)$, $\eta_{XY}:\h^*(Y)(N-q)\to \h^*(X)$
\par\noindent\textbf{Corrected:} $2(q-M)$, $\eta_{XY}:\h^*(Y)(M-q)\to \h^*(X)$
\par\noindent Integration over Y of dimension M shifts degree by2M and twist byM.
\smallskip
\item[\hypertarget{apparatus:EN-0085}{85}] Page \hyperlink{printpage:en:152}{152}\quad\texttt{S45-SOL-0626}
\par\noindent\textbf{Printed:} The proof of \ref{IV:3-4}
\par\noindent\textbf{Corrected:} The proof of \ref{IV:3-3}
\par\noindent The proof just given is that of Proposition3.3; Remark3.4 has no such proof.
\smallskip
\item[\hypertarget{apparatus:EN-0086}{86}] Page \hyperlink{printpage:en:154}{154}\quad\texttt{S45-ASTRA-R-0056}
\par\noindent\textbf{Printed:} 1
\par\noindent\textbf{Corrected:} 2
\par\noindent The printed local page number1 points to the introductory page; section1 actually begins on localpage2, printed155.
\smallskip
\item[\hypertarget{apparatus:EN-0087}{87}] Page \hyperlink{printpage:en:154}{154}\quad\texttt{S45-ASTRA-R-0057}
\par\noindent\textbf{Printed:} 12
\par\noindent\textbf{Corrected:} 13
\par\noindent The printed local page number12 points to the end of section3; section4 actually begins on localpage13, printed166.
\smallskip
\item[\hypertarget{apparatus:EN-0088}{88}] Page \hyperlink{printpage:en:155}{155}\quad\texttt{S45-SOL-0634}
\par\noindent\textbf{Printed:} \(\hom(\sL,\Hom(\sK,\dZ/n))\)
\par\noindent\textbf{Corrected:} \(\hom(\sL,\rHom(\sK,\dZ/n))\)
\par\noindent The derived tensor adjunction requires derived sheaf Hom. The scan omits R from this inner Hom, although the definition of D immediately above includes it.
\smallskip
\item[\hypertarget{apparatus:EN-0089}{89}] Page \hyperlink{printpage:en:156}{156}\quad\texttt{S45-ASTRA-R-0045}
\par\noindent\textbf{Printed:} \(\hom(j_\ast\sF,\dZ/n)=j_\ast\sF^\vee\)
\par\noindent\textbf{Corrected:} \(\Hom(j_\ast\sF,\dZ/n)=j_\ast\sF^\vee\)
\par\noindent The printed non-underlined Hom is an ordinary Hom group, but the right side is a sheaf; the sheaf Hom is required. The theorem states the correct sheaf identity.
\smallskip
\item[\hypertarget{apparatus:EN-0090}{90}] Page \hyperlink{printpage:en:156}{156}\quad\texttt{S45-SOL-0639}
\par\noindent\textbf{Printed:} \(\eR j_\ast \sF_{\bar\eta}^I\)
\par\noindent\textbf{Corrected:} \(\eR j_\ast\left(\sF_{\bar\eta}^I\right)^\vee\)
\par\noindent The bottom row is the dual triangle of the constant invariant sheaf; its rightmost term must carry the dual, as the next-page exact sequence confirms.
\smallskip
\item[\hypertarget{apparatus:EN-0091}{91}] Page \hyperlink{printpage:en:158}{158}\quad\texttt{S45-ASTRA-R-0048}
\par\noindent\textbf{Printed:} \(\rhom(f_\ast\dZ/n,\sK_Y)\)
\par\noindent\textbf{Corrected:} \(\rHom(f_\ast\dZ/n,\sK_Y)\)
\par\noindent The printed ordinary derived Hom is a complex of groups, while the equality concerns complexes of sheaves on Y; derived sheaf Hom is required.
\smallskip
\item[\hypertarget{apparatus:EN-0092}{92}] Page \hyperlink{printpage:en:158}{158}\quad\texttt{S45-SOL-0642}
\par\noindent\textbf{Printed:} \(f_\ast\dZ/n\)
\par\noindent\textbf{Corrected:} \(f_\ast\dZ/n(1)[2]\)
\par\noindent The dualizing complex K\_Y is Z/n(1)[2], so the final sheaf must carry the same twist and shift. The scan omits both.
\smallskip
\item[\hypertarget{apparatus:EN-0093}{93}] Page \hyperlink{printpage:en:162}{162}\quad\texttt{S45-SOL-0655}
\par\noindent\textbf{Printed:} \eqref{V:eq:2-2-1}
\par\noindent\textbf{Corrected:} \eqref{V:eq:3-2-1}
\par\noindent The triangle is defined by exact sequence3.2.1 immediately above, not theorem2.2.1.
\smallskip
\item[\hypertarget{apparatus:EN-0094}{94}] Page \hyperlink{printpage:en:162}{162}\quad\texttt{S45-SOL-0656}
\par\noindent\textbf{Printed:} \(\eR \pr_{2\ast}\sG\)
\par\noindent\textbf{Corrected:} \(\eR \pr_{1\ast}\sG\)
\par\noindent Both G terms use the same direct-image functor Rpr\_1*.
\smallskip
\item[\hypertarget{apparatus:EN-0095}{95}] Page \hyperlink{printpage:en:162}{162}\quad\texttt{S45-SOL-0656}
\par\noindent\textbf{Printed:} \(\eR \pr_{2\ast}\sG\)
\par\noindent\textbf{Corrected:} \(\eR \pr_{1\ast}\sG\)
\par\noindent Both G terms use the same direct-image functor Rpr\_1*.
\smallskip
\item[\hypertarget{apparatus:EN-0096}{96}] Page \hyperlink{printpage:en:163}{163}\quad\texttt{S45-SOL-0657}
\par\noindent\textbf{Printed:} \(3.2.2\)
\par\noindent\textbf{Corrected:} \(3.2.3\)
\par\noindent The print duplicates the number3.2.2 already assigned to the distinguished triangle. The next diagram is3.2.3; preserve the duplicated printed number diplomatically.
\smallskip
\item[\hypertarget{apparatus:EN-0097}{97}] Page \hyperlink{printpage:en:165}{165}\quad\texttt{S45-SOL-0664}
\par\noindent\textbf{Printed:} \(u'=c^{1,1}\)
\par\noindent\textbf{Corrected:} \(u'=-c^{1,1}\)
\par\noindent Proposition3.2 gives c\textasciicircum{}\{1,1\}=-u. The identification defining u-prime therefore gives u-prime=-c\textasciicircum{}\{1,1\}, consistent with the following trace equality.
\smallskip
\item[\hypertarget{apparatus:EN-0098}{98}] Page \hyperlink{printpage:en:173}{173}\quad\texttt{S45-SOL-0690}
\par\noindent\textbf{Printed:} \(\sF(\chi,f)\)
\par\noindent\textbf{Corrected:} \(\sF(\chi,f_0)\)
\par\noindent The morphism in the explicitly cited 1.7 d) is f\_0.
\smallskip
\item[\hypertarget{apparatus:EN-0099}{99}] Page \hyperlink{printpage:en:176}{176}\quad\texttt{S45-SOL-0696}
\par\noindent\textbf{Printed:} \(\h_c^p(X,\dQ_\ell)\)
\par\noindent\textbf{Corrected:} \(\h_c^i(X,\dQ_\ell)\)
\par\noindent The preceding isomorphism compares degrees i and i-2.
\smallskip
\item[\hypertarget{apparatus:EN-0100}{100}] Page \hyperlink{printpage:en:176}{176}\quad\texttt{S45-SOL-0697}
\par\noindent\textbf{Printed:} \(\sum (-1)^i \tr(F^\ast,\h^i(X,\dQ_\ell))\)
\par\noindent\textbf{Corrected:} \(\sum (-1)^i \tr(F^\ast,\h_c^i(X,\dQ_\ell))\)
\par\noindent The point-counting trace formula uses compact support for nonproper X.
\smallskip
\item[\hypertarget{apparatus:EN-0101}{101}] Page \hyperlink{printpage:en:178}{178}\quad\texttt{S45-SOL-0707}
\par\noindent\textbf{Printed:} $\h_c^i(X,\sF)$ and $\h^{2n-i}(X,\sF^\vee)(2n)$
\par\noindent\textbf{Corrected:} $\h_c^i(S,\sF)$ and $\h^{2n-i}(S,\sF^\vee)(n)$
\par\noindent The scheme in the hypothesis is S and the Poincaré twist in dimension n is n.
\smallskip
\item[\hypertarget{apparatus:EN-0102}{102}] Page \hyperlink{printpage:en:178}{178}\quad\texttt{S45-SOL-0710}
\par\noindent\textbf{Printed:} \(\h^0(U,\sF^\vee)(2n)\)
\par\noindent\textbf{Corrected:} \(\h^0(U,\sF^\vee)(n)\)
\par\noindent Poincaré duality in dimension n supplies twist n.
\smallskip
\item[\hypertarget{apparatus:EN-0103}{103}] Page \hyperlink{printpage:en:179}{179}\quad\texttt{S45-SOL-0711}
\par\noindent\textbf{Printed:} \((\sF_u)_{\pi_1(U,u)}(-2n)\)
\par\noindent\textbf{Corrected:} \((\sF_u)_{\pi_1(U,u)}(-n)\)
\par\noindent Dualizing the dimension-n Tate twist gives -n.
\smallskip
\item[\hypertarget{apparatus:EN-0104}{104}] Page \hyperlink{printpage:en:179}{179}\quad\texttt{S45-SOL-0719}
\par\noindent\textbf{Printed:} an open subset of $X$
\par\noindent\textbf{Corrected:} an open subset of $\bar X$
\par\noindent The open subset is X inside the proper ambient bar X.
\smallskip
\item[\hypertarget{apparatus:EN-0105}{105}] Page \hyperlink{printpage:en:180}{180}\quad\texttt{S45-SOL-0721}
\par\noindent\textbf{Printed:} (cf. \ref{VI:1-17})
\par\noindent\textbf{Corrected:} (cf. \ref{VI:1-19})
\par\noindent The comparison criterion immediately above is Proposition 1.19.
\smallskip
\item[\hypertarget{apparatus:EN-0106}{106}] Page \hyperlink{printpage:en:180}{180}\quad\texttt{S45-SOL-0722}
\par\noindent\textbf{Printed:} Apply \ref{VI:1-15}
\par\noindent\textbf{Corrected:} Apply \ref{VI:1-17}
\par\noindent The argument invokes the weight bound of Theorem 1.17.
\smallskip
\item[\hypertarget{apparatus:EN-0107}{107}] Page \hyperlink{printpage:en:181}{181}\quad\texttt{S45-SOL-0728}
\par\noindent\textbf{Printed:} $\rho(g^{-1} \gamma g)$ of $E_\gamma^n$
\par\noindent\textbf{Corrected:} $\rho(g^{-1} \gamma g)$ of $E_\lambda^n$
\par\noindent The coefficient field specified in the representation is E\_lambda; gamma labels the rational point.
\smallskip
\item[\hypertarget{apparatus:EN-0108}{108}] Page \hyperlink{printpage:en:184}{184}\quad\texttt{S45-SOL-0736}
\par\noindent\textbf{Printed:} \(\tr(F^\ast,\h^\bullet(X,\sF))\)
\par\noindent\textbf{Corrected:} \(\tr(F^\ast,\h_c^\bullet(X,\sF))\)
\par\noindent The trace formula requires compact supports without properness.
\smallskip
\item[\hypertarget{apparatus:EN-0109}{109}] Page \hyperlink{printpage:en:184}{184}\quad\texttt{S45-SOL-0737}
\par\noindent\textbf{Printed:} \(\tr(F^\ast, \h^\bullet(Y,\eR^q f_! \sF))\)
\par\noindent\textbf{Corrected:} \(\tr(F^\ast,\h_c^\bullet(Y,\eR^q f_! \sF))\)
\par\noindent The trace formula on Y is a compact-support trace formula.
\smallskip
\item[\hypertarget{apparatus:EN-0110}{110}] Page \hyperlink{printpage:en:185}{185}\quad\texttt{S45-SOL-0738}
\par\noindent\textbf{Printed:} \(\tr(F^\ast,\h^\bullet(Y,\eR^\bullet f_! \sF))\)
\par\noindent\textbf{Corrected:} \(\tr(F^\ast,\h_c^\bullet(Y,\eR^\bullet f_! \sF))\)
\par\noindent Alternating the compact-support Leray sequence preserves compact supports.
\smallskip
\item[\hypertarget{apparatus:EN-0111}{111}] Page \hyperlink{printpage:en:187}{187}\quad\texttt{S45-SOL-0748}
\par\noindent\textbf{Printed:} \((\chi^{-1}(g)\)
\par\noindent\textbf{Corrected:} \(\chi^{-1}(g)\)
\par\noindent Delete the unmatched opening parenthesis.
\smallskip
\item[\hypertarget{apparatus:EN-0112}{112}] Page \hyperlink{printpage:en:187}{187}\quad\texttt{S45-SOL-0749}
\par\noindent\textbf{Printed:} \(\h^\bullet(G,\sF(\chi))\)
\par\noindent\textbf{Corrected:} \(\h_c^\bullet(G,\sF(\chi))\)
\par\noindent rho\_H was defined on compactly supported cohomology.
\smallskip
\item[\hypertarget{apparatus:EN-0113}{113}] Page \hyperlink{printpage:en:188}{188}\quad\texttt{S45-SOL-0750}
\par\noindent\textbf{Printed:} \(\h^p(X,\sF)\)
\par\noindent\textbf{Corrected:} \(\h_c^p(X,\sF)\)
\par\noindent The fibre of R\textasciicircum{}p pr\_\{1!\} has compact supports.
\smallskip
\item[\hypertarget{apparatus:EN-0114}{114}] Page \hyperlink{printpage:en:189}{189}\quad\texttt{S45-SOL-0757}
\par\noindent\textbf{Printed:} \(\h^{2-i}(\bar X,j_\ast(\sF^\vee(1))\)
\par\noindent\textbf{Corrected:} \(\h^{2-i}(\bar X,j_\ast(\sF^\vee(1)))\)
\par\noindent Balance the parentheses of the cohomology group.
\smallskip
\item[\hypertarget{apparatus:EN-0115}{115}] Page \hyperlink{printpage:en:190}{190}\quad\texttt{S45-SOL-0760}
\par\noindent\textbf{Printed:} \(\psi(a^p-a)=0\)
\par\noindent\textbf{Corrected:} \(\psi(a^p-a)=1\)
\par\noindent The defining exponential has value 1 when the trace is zero.
\smallskip
\item[\hypertarget{apparatus:EN-0116}{116}] Page \hyperlink{printpage:en:190}{190}\quad\texttt{S45-SOL-0763}
\par\noindent\textbf{Printed:} This sum is zero for $f$ of the form $g^p-g$.
\par\noindent\textbf{Corrected:} For $f$ of the form $g^p-g$, this sum equals $\#\{x\in X_0(\dF_q)\mid f(x)\ne\infty\}$.
\par\noindent Every summand in this case is 1.
\smallskip
\item[\hypertarget{apparatus:EN-0117}{117}] Page \hyperlink{printpage:en:190}{190}\quad\texttt{S45-SOL-0765}
\par\noindent\textbf{Printed:} for every closed point $s$
\par\noindent\textbf{Corrected:} for every closed point $x$
\par\noindent Subsequent notation uniformly uses x.
\smallskip
\item[\hypertarget{apparatus:EN-0118}{118}] Page \hyperlink{printpage:en:191}{191}\quad\texttt{S45-SOL-0771}
\par\noindent\textbf{Printed:} that of the $\dF_q$-torsor
\par\noindent\textbf{Corrected:} that of the $\dF_p$-torsor
\par\noindent The kernel of the Artin–Schreier polynomial T\textasciicircum{}p-T is F\_p.
\smallskip
\item[\hypertarget{apparatus:EN-0119}{119}] Page \hyperlink{printpage:en:192}{192}\quad\texttt{S45-SOL-0773}
\par\noindent\textbf{Printed:} $\dF_q$-torsor over $\spec(\dF_q)$
\par\noindent\textbf{Corrected:} $\dF_p$-torsor over $\spec(\dF_q)$
\par\noindent The structure group is F\_p whereas the base remains Spec(F\_q).
\smallskip
\item[\hypertarget{apparatus:EN-0120}{120}] Page \hyperlink{printpage:en:192}{192}\quad\texttt{S45-SOL-0780}
\par\noindent\textbf{Printed:} perfect duality with values in $E_\lambda(1)$
\par\noindent\textbf{Corrected:} perfect duality with values in $E_\lambda(-1)$
\par\noindent Poincaré duality pairs degree-one groups into H\textasciicircum{}2(E\_lambda)=E\_lambda(-1).
\smallskip
\item[\hypertarget{apparatus:EN-0121}{121}] Page \hyperlink{printpage:en:192}{192}\quad\texttt{S45-SOL-0781}
\par\noindent\textbf{Printed:} $\sigma^\ast$ acts trivially on $E_\lambda(1)$
\par\noindent\textbf{Corrected:} $\sigma^\ast$ acts trivially on $E_\lambda(-1)$
\par\noindent Maintain the corrected Poincaré-duality target.
\smallskip
\item[\hypertarget{apparatus:EN-0122}{122}] Page \hyperlink{printpage:en:193}{193}\quad\texttt{S45-SOL-0786}
\par\noindent\textbf{Printed:} $r\geqslant i$
\par\noindent\textbf{Corrected:} $r\geqslant 1$
\par\noindent The lower bound in the induction is the integer 1.
\smallskip
\item[\hypertarget{apparatus:EN-0123}{123}] Page \hyperlink{printpage:en:194}{194}\quad\texttt{S45-SOL-0792}
\par\noindent\textbf{Printed:} {\footnotesize\ttfamily\raggedright \textbackslash{}\allowbreak{}sum\_\allowbreak{}\{\allowbreak{}x\_\allowbreak{}2,\allowbreak{}\textbackslash{}\allowbreak{}dots,\allowbreak{}x\_\allowbreak{}n\}\allowbreak{} \textbackslash{}\allowbreak{}sum\_\allowbreak{}\{\allowbreak{}x\_\allowbreak{}1\}\allowbreak{} \textbackslash{}\allowbreak{}psi P(\allowbreak{}x\_\allowbreak{}1,\allowbreak{}\textbackslash{}\allowbreak{}dots,\allowbreak{}x\_\allowbreak{}n\textbackslash{}\allowbreak{}right|}
\par\noindent\textbf{Corrected:} {\footnotesize\ttfamily\raggedright \textbackslash{}\allowbreak{}sum\_\allowbreak{}\{\allowbreak{}x\_\allowbreak{}2,\allowbreak{}\textbackslash{}\allowbreak{}dots,\allowbreak{}x\_\allowbreak{}n\}\allowbreak{} \textbackslash{}\allowbreak{}sum\_\allowbreak{}\{\allowbreak{}x\_\allowbreak{}1\}\allowbreak{} \textbackslash{}\allowbreak{}psi P(\allowbreak{}x\_\allowbreak{}1,\allowbreak{}\textbackslash{}\allowbreak{}dots,\allowbreak{}x\_\allowbreak{}n)\allowbreak{}\textbackslash{}\allowbreak{}right|}
\par\noindent Close the argument of P before the absolute-value delimiter.
\smallskip
\item[\hypertarget{apparatus:EN-0124}{124}] Page \hyperlink{printpage:en:195}{195}\quad\texttt{S45-SOL-0794}
\par\noindent\textbf{Printed:} {\footnotesize\ttfamily\raggedright \textbackslash{}\allowbreak{}left.}
\par\noindent\textbf{Corrected:} {\footnotesize\ttfamily\raggedright \textbackslash{}\allowbreak{}left|}
\par\noindent Supply the missing opening absolute-value delimiter.
\smallskip
\item[\hypertarget{apparatus:EN-0125}{125}] Page \hyperlink{printpage:en:195}{195}\quad\texttt{S45-SOL-0795}
\par\noindent\textbf{Printed:} \(\sum_{S=\lambda}\)
\par\noindent\textbf{Corrected:} \(\sum_{S_1=\lambda}\)
\par\noindent The linear form in the argument is S\_1.
\smallskip
\item[\hypertarget{apparatus:EN-0126}{126}] Page \hyperlink{printpage:en:195}{195}\quad\texttt{S45-SOL-0796}
\par\noindent\textbf{Printed:} \(\sum_{S(X_2,\dots,X_n)\ne\lambda}\)
\par\noindent\textbf{Corrected:} \(\sum_{S_1(X_2,\dots,X_n)\ne\lambda}\)
\par\noindent The same linear form is S\_1.
\smallskip
\item[\hypertarget{apparatus:EN-0127}{127}] Page \hyperlink{printpage:en:195}{195}\quad\texttt{S45-SOL-0797}
\par\noindent\textbf{Printed:} \(< (d-2) q^{n-1/2} + q^{n-1}\)
\par\noindent\textbf{Corrected:} \(\leqslant (d-2) q^{n-1/2} + q^{n-1}\)
\par\noindent Equality occurs at this intermediate step when d=2.
\smallskip
\item[\hypertarget{apparatus:EN-0128}{128}] Page \hyperlink{printpage:en:196}{196}\quad\texttt{S45-SOL-0799}
\par\noindent\textbf{Printed:} \(\tr(F^\ast,\h_c^\ast(\dG_m,\sF(\psi \chi^{-1}))\)
\par\noindent\textbf{Corrected:} \(\tr(F^\ast,\h_c^\ast(\dG_m,\sF(\psi \chi^{-1})))\)
\par\noindent Close the trace expression.
\smallskip
\item[\hypertarget{apparatus:EN-0129}{129}] Page \hyperlink{printpage:en:197}{197}\quad\texttt{S45-SOL-0805}
\par\noindent\textbf{Printed:} \(|\tau(\chi,\psi)|=q\)
\par\noindent\textbf{Corrected:} \(|\tau(\chi,\psi)|=q^{1/2}\)
\par\noindent The product with the conjugate is the square of the absolute value.
\smallskip
\item[\hypertarget{apparatus:EN-0130}{130}] Page \hyperlink{printpage:en:198}{198}\quad\texttt{S45-SOL-0811}
\par\noindent\textbf{Printed:} $Q$ $P$ over
\par\noindent\textbf{Corrected:} $Q$ and $P$ over
\par\noindent Restore the missing conjunction joining the two projective spaces.
\smallskip
\item[\hypertarget{apparatus:EN-0131}{131}] Page \hyperlink{printpage:en:200}{200}\quad\texttt{S45-SOL-0819}
\par\noindent\textbf{Printed:} \(\gal(\bar\dF/F)\)
\par\noindent\textbf{Corrected:} \(\gal(\dF/\dF_q)\)
\par\noindent Arithmetic Galois action over the finite base field.
\smallskip
\item[\hypertarget{apparatus:EN-0132}{132}] Page \hyperlink{printpage:en:200}{200}\quad\texttt{S45-SOL-0820}
\par\noindent\textbf{Printed:} an embedding $\omega$ of $k$ into $\dF_q$
\par\noindent\textbf{Corrected:} an embedding $\omega$ of $k$ into $\dF$
\par\noindent A general finite extension k embeds into the algebraic closure, not the base field.
\smallskip
\item[\hypertarget{apparatus:EN-0133}{133}] Page \hyperlink{printpage:en:202}{202}\quad\texttt{S45-SOL-0828}
\par\noindent\textbf{Printed:} \(\tr_{\dF_q/k}\)
\par\noindent\textbf{Corrected:} \(\tr_{k/\dF_q}\)
\par\noindent Reverse the incorrectly ordered source and target of the trace.
\smallskip
\item[\hypertarget{apparatus:EN-0134}{134}] Page \hyperlink{printpage:en:202}{202}\quad\texttt{S45-SOL-0830}
\par\noindent\textbf{Printed:} \(N^\ast\sF(\chi) = \bigotimes_i \pr_i^\ast\sF(\psi)\)
\par\noindent\textbf{Corrected:} \(N^\ast\sF(\chi) = \bigotimes_i \pr_i^\ast\sF(\chi)\)
\par\noindent The norm pullback factors into copies of the multiplicative-character sheaf.
\smallskip
\item[\hypertarget{apparatus:EN-0135}{135}] Page \hyperlink{printpage:en:202}{202}\quad\texttt{S45-SOL-0832}
\par\noindent\textbf{Printed:} the cardinality $n$ of $I$
\par\noindent\textbf{Corrected:} the cardinality $N$ of $I$
\par\noindent The cardinality of I is N.
\smallskip
\item[\hypertarget{apparatus:EN-0136}{136}] Page \hyperlink{printpage:en:202}{202}\quad\texttt{S45-SOL-0834}
\par\noindent\textbf{Printed:} \(\h_c^1(\dG_m,\sF(\psi\chi^{-1})) \otimes_\dZ\)
\par\noindent\textbf{Corrected:} \(\h_c^1(\dG_m,\sF(\psi\chi^{-1}))^{\otimes N} \otimes_\dZ\)
\par\noindent Künneth requires N copies of the degree-one cohomology.
\smallskip
\item[\hypertarget{apparatus:EN-0137}{137}] Page \hyperlink{printpage:en:202}{202}\quad\texttt{S45-SOL-0833}
\par\noindent\textbf{Printed:} {\footnotesize\ttfamily\raggedright \textbackslash{}\allowbreak{}h\_\allowbreak{}c\textasciicircum{}\allowbreak{}N\textbackslash{}\allowbreak{}left(\allowbreak{}V,\allowbreak{}}
\par\noindent\textbf{Corrected:} {\footnotesize\ttfamily\raggedright \textbackslash{}\allowbreak{}h\_\allowbreak{}c\textasciicircum{}\allowbreak{}N\textbackslash{}\allowbreak{}left(\allowbreak{}V\textasciicircum{}\allowbreak{}\textbackslash{}\allowbreak{}ast,\allowbreak{}}
\par\noindent The character sheaf is on the unit torus.
\smallskip
\item[\hypertarget{apparatus:EN-0138}{138}] Page \hyperlink{printpage:en:203}{203}\quad\texttt{S45-SOL-0839}
\par\noindent\textbf{Printed:} \((k^\ast)^{N+1}\)
\par\noindent\textbf{Corrected:} \((\dF_q^\ast)^{N+1}\)
\par\noindent The coordinate tuple is printed as an element of k\textasciicircum{}\{* N+1\}, but each coordinate belongs to F\_q\textasciicircum{}* because k is the product of N+1 copies of F\_q. Correct the coordinate field.
\smallskip
\item[\hypertarget{apparatus:EN-0139}{139}] Page \hyperlink{printpage:en:203}{203}\quad\texttt{S45-SOL-0840}
\par\noindent\textbf{Printed:} \(n\)
\par\noindent\textbf{Corrected:} \(N\)
\par\noindent The upper limit in the trace-zero condition is printed as lowercase n. The coordinates are indexed from 1 to N, so the bound is N.
\smallskip
\item[\hypertarget{apparatus:EN-0140}{140}] Page \hyperlink{printpage:en:203}{203}\quad\texttt{S45-SOL-0841}
\par\noindent\textbf{Printed:} \(\chi_n^{-1}(x_n)\)
\par\noindent\textbf{Corrected:} \(\chi_N^{-1}(x_N)\)
\par\noindent The final factor in (4.14.2) visibly has lowercase n in both chi\_n and x\_n, although the summation uses N coordinates. Correct both indices to N.
\smallskip
\item[\hypertarget{apparatus:EN-0141}{141}] Page \hyperlink{printpage:en:204}{204}\quad\texttt{S45-SOL-0843}
\par\noindent\textbf{Printed:} \(q\sum_{x\in k^\ast/\dF_q^\ast} \chi(x)^{-1}\)
\par\noindent\textbf{Corrected:} \(q\sum_{\substack{x\in k^\ast/\dF_q^\ast\\\tr(x)=0}} \chi(x)^{-1}\)
\par\noindent The first term comes from the trace-zero indicator.
\smallskip
\item[\hypertarget{apparatus:EN-0142}{142}] Page \hyperlink{printpage:en:205}{205}\quad\texttt{S45-SOL-0847}
\par\noindent\textbf{Printed:} \(\sK_m((\chi_i)_{i\in I})\)
\par\noindent\textbf{Corrected:} \(\sK_n((\chi_i)_{i\in I})\)
\par\noindent The Kummer exponent of the given characters is n.
\smallskip
\item[\hypertarget{apparatus:EN-0143}{143}] Page \hyperlink{printpage:en:205}{205}\quad\texttt{S45-SOL-0847}
\par\noindent\textbf{Printed:} \(\sK_m((\chi_i)_{i\in I})\)
\par\noindent\textbf{Corrected:} \(\sK_n((\chi_i)_{i\in I})\)
\par\noindent The Kummer exponent of the given characters is n.
\smallskip
\item[\hypertarget{apparatus:EN-0144}{144}] Page \hyperlink{printpage:en:205}{205}\quad\texttt{S45-SOL-0849}
\par\noindent\textbf{Printed:} $T_0 V_0$ vanishes
\par\noindent\textbf{Corrected:} $T_0 v_0$ vanishes
\par\noindent Composition with the blow-up morphism v\_0.
\smallskip
\item[\hypertarget{apparatus:EN-0145}{145}] Page \hyperlink{printpage:en:205}{205}\quad\texttt{S45-SOL-0850}
\par\noindent\textbf{Printed:} on $\pi^{-1}(Q_0)$
\par\noindent\textbf{Corrected:} on $\bar\pi^{-1}(Q_0)$
\par\noindent Use the projection of the full line bundle.
\smallskip
\item[\hypertarget{apparatus:EN-0146}{146}] Page \hyperlink{printpage:en:205}{205}\quad\texttt{S45-SOL-0851}
\par\noindent\textbf{Printed:} $D=\pi^{-1}(x)$
\par\noindent\textbf{Corrected:} $D=\bar\pi^{-1}(x)$
\par\noindent The full line is the fibre of bar pi.
\smallskip
\item[\hypertarget{apparatus:EN-0147}{147}] Page \hyperlink{printpage:en:206}{206}\quad\texttt{S45-SOL-0856}
\par\noindent\textbf{Printed:} \(N-1\)
\par\noindent\textbf{Corrected:} \(N-2\)
\par\noindent With the revised rank N of M, the projective trace-zero locus has dimension N-2.
\smallskip
\item[\hypertarget{apparatus:EN-0148}{148}] Page \hyperlink{printpage:en:206}{206}\quad\texttt{S45-SOL-0856}
\par\noindent\textbf{Printed:} \(N-1\)
\par\noindent\textbf{Corrected:} \(N-2\)
\par\noindent With the revised rank N of M, the projective trace-zero locus has dimension N-2.
\smallskip
\item[\hypertarget{apparatus:EN-0149}{149}] Page \hyperlink{printpage:en:206}{206}\quad\texttt{S45-SOL-0856}
\par\noindent\textbf{Printed:} \(N-1\)
\par\noindent\textbf{Corrected:} \(N-2\)
\par\noindent With the revised rank N of M, the projective trace-zero locus has dimension N-2.
\smallskip
\item[\hypertarget{apparatus:EN-0150}{150}] Page \hyperlink{printpage:en:208}{208}\quad\texttt{S45-SOL-0859}
\par\noindent\textbf{Printed:} and $K$ a field
\par\noindent\textbf{Corrected:} and $k$ a field
\par\noindent The field is denoted k throughout the subsequent definitions.
\smallskip
\item[\hypertarget{apparatus:EN-0151}{151}] Page \hyperlink{printpage:en:210}{210}\quad\texttt{S45-SOL-0879}
\par\noindent\textbf{Printed:} \(\chi_1|F_1=(N_{F_1/\dQ})^N\)
\par\noindent\textbf{Corrected:} \(\chi_1|F_1'=(N_{F_1'/\dQ})^N\)
\par\noindent The restriction is to the totally real subfield, not to the entire CM field.
\smallskip
\item[\hypertarget{apparatus:EN-0152}{152}] Page \hyperlink{printpage:en:210}{210}\quad\texttt{S45-SOL-0882}
\par\noindent\textbf{Printed:} $k$ has torsion
\par\noindent\textbf{Corrected:} $k$ is torsion-free
\par\noindent The fractional-ideal group is torsion-free, as used by the argument.
\smallskip
\item[\hypertarget{apparatus:EN-0153}{153}] Page \hyperlink{printpage:en:213}{213}\quad\texttt{S45-SOL-0889}
\par\noindent\textbf{Printed:} K\"ummer
\par\noindent\textbf{Corrected:} Kummer
\par\noindent Correct the proper name.
\smallskip
\item[\hypertarget{apparatus:EN-0154}{154}] Page \hyperlink{printpage:en:214}{214}\quad\texttt{S45-SOL-0898}
\par\noindent\textbf{Printed:} embeddings of $e_\beta$ into $f$
\par\noindent\textbf{Corrected:} embeddings of $e_\beta$ into $\bar f$
\par\noindent The embeddings are into the algebraic closure of the residue field.
\smallskip
\item[\hypertarget{apparatus:EN-0155}{155}] Page \hyperlink{printpage:en:215}{215}\quad\texttt{S45-SOL-0899}
\par\noindent\textbf{Printed:} \(\lambda_{d(\beta)}\)
\par\noindent\textbf{Corrected:} \(\lambda_{\alpha(\beta)}\)
\par\noindent The character is indexed by the orbit alpha(beta).
\smallskip
\item[\hypertarget{apparatus:EN-0156}{156}] Page \hyperlink{printpage:en:215}{215}\quad\texttt{S45-SOL-0903}
\par\noindent\textbf{Printed:} \((\chi_\beta\circ \sigma_i^{-1})_{i\in I}\)
\par\noindent\textbf{Corrected:} \((\chi_\beta\circ \sigma_i^{-1})_{i\in\beta}\)
\par\noindent The component torus uses the embeddings in beta.
\smallskip
\item[\hypertarget{apparatus:EN-0157}{157}] Page \hyperlink{printpage:en:215}{215}\quad\texttt{S45-SOL-0904}
\par\noindent\textbf{Printed:} \(\sK_n((_n\bar\lambda_i)_{i\in \beta}) = \sF(\chi)\)
\par\noindent\textbf{Corrected:} \(\sK_n((_n\bar\lambda_i)_{i\in I}) = \sF(\chi)\)
\par\noindent The global torus uses all embeddings I.
\smallskip
\item[\hypertarget{apparatus:EN-0158}{158}] Page \hyperlink{printpage:en:215}{215}\quad\texttt{S45-SOL-0906}
\par\noindent\textbf{Printed:} \((\chi_\beta)_{\beta\in B}\)
\par\noindent\textbf{Corrected:} \((\chi_\beta)_{\beta\in D}\)
\par\noindent The orbit set is D, not the undefined B.
\smallskip
\item[\hypertarget{apparatus:EN-0159}{159}] Page \hyperlink{printpage:en:216}{216}\quad\texttt{S45-SOL-0908}
\par\noindent\textbf{Printed:} a set of density $0$
\par\noindent\textbf{Corrected:} a set of density $1$
\par\noindent Degree-one prime ideals have density one among prime ideals of the number field.
\smallskip
\item[\hypertarget{apparatus:EN-0160}{160}] Page \hyperlink{printpage:en:217}{217}\quad\texttt{S45-ASTRA-S-0037}
\par\noindent\textbf{Printed:} $a\in A_j\subset (\dZ/n)^\ast$
\par\noindent\textbf{Corrected:} $a\in A_j\subset\dZ/n$
\par\noindent The orbit family is in all of Z/n, as defined in 6.9(a).
\smallskip
\item[\hypertarget{apparatus:EN-0161}{161}] Page \hyperlink{printpage:en:220}{220}\quad\texttt{S45-SOL-0934}
\par\noindent\textbf{Printed:} \(\kloos_{\dF_{q^n},a}\)
\par\noindent\textbf{Corrected:} \(\kloos_{\dF_q^n,a}\)
\par\noindent The original n-variable sum uses the split étale algebra F\_q\textasciicircum{}n.
\smallskip
\item[\hypertarget{apparatus:EN-0162}{162}] Page \hyperlink{printpage:en:226}{226}\quad\texttt{S45-SOL-0951}
\par\noindent\textbf{Printed:} eigenvalues of $F$ is therefore $q^{\frac{1-n}{2}}$
\par\noindent\textbf{Corrected:} eigenvalues of $F$ is therefore $q^{\frac{n-1}{2}}$
\par\noindent Undo the normalization factor of F to recover its eigenvalue.
\smallskip
\item[\hypertarget{apparatus:EN-0163}{163}] Page \hyperlink{printpage:en:227}{227}\quad\texttt{S45-SOL-0959}
\par\noindent\textbf{Printed:} functions on $X$ invariant under $\tau$
\par\noindent\textbf{Corrected:} functions on $V_a^1\times V_a^1$ invariant under $\tau$
\par\noindent Invariant functions are on the space before taking the quotient; also restore plural agreement.
\smallskip
\item[\hypertarget{apparatus:EN-0164}{164}] Page \hyperlink{printpage:en:228}{228}\quad\texttt{S45-SOL-0966}
\par\noindent\textbf{Printed:} The symmetric group $\mathfrak{S}_m$
\par\noindent\textbf{Corrected:} The symmetric group $\mathfrak{S}_n$
\par\noindent The group permutes n coordinates.
\smallskip
\item[\hypertarget{apparatus:EN-0165}{165}] Page \hyperlink{printpage:en:229}{229}\quad\texttt{S45-SOL-0973}
\par\noindent\textbf{Printed:} \(\h_0^\bullet(\dG_a,\sF(\psi))\)
\par\noindent\textbf{Corrected:} \(\h_c^\bullet(\dG_a,\sF(\psi))\)
\par\noindent Use compactly supported cohomology in the Künneth vanishing.
\smallskip
\item[\hypertarget{apparatus:EN-0166}{166}] Page \hyperlink{printpage:en:230}{230}\quad\texttt{S45-SOL-0978}
\par\noindent\textbf{Printed:} Weil group $W$
\par\noindent\textbf{Corrected:} Weyl group $W$
\par\noindent The group acting on the maximal torus is the Weyl group.
\smallskip
\item[\hypertarget{apparatus:EN-0167}{167}] Page \hyperlink{printpage:en:232}{232}\quad\texttt{S45-ASTRA-S-0049}
\par\noindent\textbf{Printed:} Weil
\par\noindent\textbf{Corrected:} Weyl
\par\noindent Correct the proper name of the reflection groups.
\smallskip
\item[\hypertarget{apparatus:EN-0168}{168}] Page \hyperlink{printpage:en:234}{234}\quad\texttt{S45-SOL-0981}
\par\noindent\textbf{Printed:} \(E_1^{p q}\)
\par\noindent\textbf{Corrected:} \(E_2^{p q}\)
\par\noindent The first display on p.234 visibly has E\_1. The displayed derived-functor cohomology terms constitute the E\_2 page, as in the Leray display immediately below; retain 1 diplomatically and correct to 2.
\smallskip
\item[\hypertarget{apparatus:EN-0169}{169}] Page \hyperlink{printpage:en:234}{234}\quad\texttt{S45-SOL-0989}
\par\noindent\textbf{Printed:} \(\eR g^! K(n)[2 n]\)
\par\noindent\textbf{Corrected:} \(\eR g^! K \iso g^\ast K(n)[2 n]\)
\par\noindent The bottom line contains Rg\textasciicircum{}!K(n)[2n] without an equality or pullback. Smooth duality and the stated relative dimension require Rg\textasciicircum{}!K ≃ g\textasciicircum{}*K(n)[2n].
\smallskip
\item[\hypertarget{apparatus:EN-0170}{170}] Page \hyperlink{printpage:en:235}{235}\quad\texttt{S45-SOL-0990}
\par\noindent\textbf{Printed:} \(K\to \eR j_\ast j^\ast K\)
\par\noindent\textbf{Corrected:} \(L\to \eR j_\ast j^\ast L\)
\par\noindent At the top of p.235 the variable L is declared, but the displayed arrow contains K twice. The natural map for the same object L requires L in both positions.
\smallskip
\item[\hypertarget{apparatus:EN-0171}{171}] Page \hyperlink{printpage:en:235}{235}\quad\texttt{S45-SOL-0992}
\par\noindent\textbf{Printed:} \(\eR j_\ast \Hom(\sF_1,\eR j^! \sG)\)
\par\noindent\textbf{Corrected:} \(\eR j_\ast \rHom(\sF_1,\eR j^! \sG)\)
\par\noindent The right side of the printed derived Hom identity lacks R before internal Hom. Its left side and the subsequent derived assertion require the derived internal Hom; the correction is substantive, not a choice of sheaf typography.
\smallskip
\item[\hypertarget{apparatus:EN-0172}{172}] Page \hyperlink{printpage:en:236}{236}\quad\texttt{S45-SOL-0996}
\par\noindent\textbf{Printed:} Let $F:X\to Y$ be a morphism of $S$-schemes of finite type
\par\noindent\textbf{Corrected:} Let $f:X\to Y$ be a morphism of $S$-schemes of finite type
\par\noindent Theorem1.9 starts with capital F, while both its conclusions use lowercase f. The capital is a misprint; preserve it only diplomatically.
\smallskip
\item[\hypertarget{apparatus:EN-0173}{173}] Page \hyperlink{printpage:en:239}{239}\quad\texttt{S45-SOL-0999}
\par\noindent\textbf{Printed:} \(Y\subset \dA_S^n\)
\par\noindent\textbf{Corrected:} \(Y\subset \dA_S^N\)
\par\noindent At the bottom of239 the affine embedding uses n, already fixed as the relative-dimension bound in(*\_n). An arbitrary affine scheme of dimension at most n need not embed in affine n-space; distinguish the embedding dimension N.
\smallskip
\item[\hypertarget{apparatus:EN-0174}{174}] Page \hyperlink{printpage:en:243}{243}\quad\texttt{S45-SOL-1015}
\par\noindent\textbf{Printed:} \ref{VII:2-13} is always equal to $S$.
\par\noindent\textbf{Corrected:} \ref{VII:2-13} is always equal to $S$).
\par\noindent The last sentence opens a parenthesis after limite and never closes it before the final full stop. Supply the missing parenthesis only in corrected forms.
\smallskip
\item[\hypertarget{apparatus:EN-0175}{175}] Page \hyperlink{printpage:en:246}{246}\quad\texttt{S45-SOL-1031}
\par\noindent\textbf{Printed:} \(\h^p(X_s,\eR^i\Psi_\eta(\sF))\)
\par\noindent\textbf{Corrected:} \(\h^p(X_s,\eR^q\Psi_\eta(\sF))\)
\par\noindent The bottom spectral sequence prints i as the nearby-cycle degree, although the term is E\_2\textasciicircum{}\{pq\} and abuts in degree p+q. The nearby-cycle degree must be q; subsequent G\textasciicircum{}q and H\textasciicircum{}q confirm it.
\smallskip
\item[\hypertarget{apparatus:EN-0176}{176}] Page \hyperlink{printpage:en:248}{248}\quad\texttt{S45-SOL-1035}
\par\noindent\textbf{Printed:} $\eR\Gamma(X_s,\Delta)=0$
\par\noindent\textbf{Corrected:} $\eR\Gamma(X_{s'}',\Delta)=0$
\par\noindent Assertion(B) visibly uses X\_s. Its Δ was defined from the base-change morphism on X′\_\{s′\}; paragraph3.9 itself computes RΓ(X′\_\{s′\},Δ). Correct the space accordingly.
\smallskip
\item[\hypertarget{apparatus:EN-0177}{177}] Page \hyperlink{printpage:en:250}{250}\quad\texttt{S45-SOL-1047}
\par\noindent\textbf{Printed:} $K\in\ob\eD_\text{ctf}^b(S,A)$, put
\par\noindent\textbf{Corrected:} $K\in\ob\eD_\text{ctf}^b(X,A)$, put
\par\noindent Paragraph4.2 assigns K to D\_ctf\textasciicircum{}b(S,A), although K is dualized against K\_X and theorem4.3 concerns complexes on X. The printed S is a base-space error; use X in corrected forms.
\smallskip
\item[\hypertarget{apparatus:EN-0178}{178}] Page \hyperlink{printpage:en:253}{253}\quad\texttt{S45-SOL-1056}
\par\noindent\textbf{Printed:} \(X_\bullet=X\times_S \bar S_\bullet\)
\par\noindent\textbf{Corrected:} \(\bar X_\bullet=X\times_S \bar S_\bullet\)
\par\noindent The printed definition uses unbarred X\_bullet, but the very next projection is from barred X\_bullet. Supply the bar at the definition; retain the printed fiber product X×\_S barred S\_bullet, which is already correct.
\smallskip
\item[\hypertarget{apparatus:EN-0179}{179}] Page \hyperlink{printpage:en:253}{253}\quad\texttt{S45-ASTRA-V-0004}
\par\noindent\textbf{Printed:} Put
$G|X'=G'$. The bicomplex ${g_\bullet}_\ast g_\bullet^\ast G'$, whose columns
are the ${g_n}_\ast g_n^\ast G'$, is a resolution of $G'$
\par\noindent\textbf{Corrected:} Put
$E|X'=E'$. The bicomplex ${g_\bullet}_\ast g_\bullet^\ast E'$, whose columns
are the ${g_n}_\ast g_n^\ast E'$, is a resolution of $E'$
\par\noindent The proposition fixes E and constructs the base-change cone from E. No object G is introduced before the resolution, which must resolve E′ to compute that cone. The scan clearly prints G in all five places.
\smallskip
\item[\hypertarget{apparatus:EN-0180}{180}] Page \hyperlink{printpage:en:255}{255}\quad\texttt{S45-SOL-1062}
\par\noindent\textbf{Printed:} $\widetilde D$ of $D$ in $Y$
\par\noindent\textbf{Corrected:} $\widetilde D$ of $D$ in $\widetilde Y$
\par\noindent The proof prints the inverse image D-tilde in Y, although the newly constructed covering is Y-tilde→Y. The inverse-image divisor belongs on Y-tilde.
\smallskip
\item[\hypertarget{apparatus:EN-0181}{181}] Page \hyperlink{printpage:en:255}{255}\quad\texttt{S45-SOL-1063}
\par\noindent\textbf{Printed:} \(f_i:Y\setminus D_i \to D\)
\par\noindent\textbf{Corrected:} \(f_i:Y\setminus D_i \to Y\)
\par\noindent The printed inclusion f\_i:Y−D\_i→D has target D. Its domain is an open subset of Y and the stated canonical inclusion must end at Y; the following direct-image formulas also live on Y.
\smallskip
\item[\hypertarget{apparatus:EN-0182}{182}] Page \hyperlink{printpage:en:257}{257}\quad\texttt{S45-SOL-1070}
\par\noindent\textbf{Printed:} \(\eR\Gamma(\bar X,E|\bar X) (=\eR\bar f_\ast (E|\bar X)_s \to\)
\par\noindent\textbf{Corrected:} \(\eR\Gamma(\bar X,E|\bar X) (=\eR\bar f_\ast (E|\bar X)_s) \to\)
\par\noindent The first equation on257 visibly opens a parenthesis before the equality and leaves it unclosed. Miller removes that opening mark. Restore it diplomatically; close the parenthetical equality before the arrow in corrected forms.
\smallskip
\item[\hypertarget{apparatus:EN-0183}{183}] Page \hyperlink{printpage:en:257}{257}\quad\texttt{S45-SOL-1076}
\par\noindent\textbf{Printed:} \(\spe(u')\spe(u'')\)
\par\noindent\textbf{Corrected:} \(\spe(u')\spe(u)\)
\par\noindent The specialization formula prints sp(u′)sp(u″), although the arrows just defined are u′ and u only. Functorial contravariant composition requires sp(u′)sp(u).
\smallskip
\item[\hypertarget{apparatus:EN-0184}{184}] Page \hyperlink{printpage:en:259}{259}\quad\texttt{S45-SOL-1085}
\par\noindent\textbf{Printed:} \(\eR\Gamma(Y(y)_t,(\eR {f_{(x)}}_\ast (E|X(x))|Y(y)_t)\)
\par\noindent\textbf{Corrected:} \(\eR\Gamma\left(Y(y)_t,\left(\eR {f_{(x)}}_\ast (E|X(x))\right)|Y(y)_t\right)\)
\par\noindent The right-hand cohomology term visibly opens a parenthesis before Rf\_(x)* but does not close that grouping before restriction. Supply the missing closing parenthesis around the complex in corrected forms; preserve the printed defect diplomatically.
\smallskip
\item[\hypertarget{apparatus:EN-0185}{185}] Page \hyperlink{printpage:en:260}{260}\quad\texttt{S45-SOL-1086}
\par\noindent\textbf{Printed:} \(\eR f'_\ast(E X')\)
\par\noindent\textbf{Corrected:} \(\eR f'_\ast(E|X')\)
\par\noindent The first proof paragraph prints Rf′\_*(E X′), omitting the restriction bar between E and X′. The argument and surrounding notation require E|X′; Miller has silently supplied it.
\smallskip
\item[\hypertarget{apparatus:EN-0186}{186}] Page \hyperlink{printpage:en:260}{260}\quad\texttt{S45-SOL-1087}
\par\noindent\textbf{Printed:} image of
$s$ (resp.\ $t$) in $S$
\par\noindent\textbf{Corrected:} image of
$s'$ (resp.\ $t'$) in $S$
\par\noindent The scan calls s and t the images of s and t, without primes. The source points just defined lie on S′ and are s′,t′; the missing primes must be supplied in the corrected text.
\smallskip
\item[\hypertarget{apparatus:EN-0187}{187}] Page \hyperlink{printpage:en:260}{260}\quad\texttt{S45-SOL-1089}
\par\noindent\textbf{Printed:} \(\eR\Gamma(X(x)_t,(E|X(x)_t)\lotimes_A M\)
\par\noindent\textbf{Corrected:} \(\eR\Gamma(X(x)_t,(E|X(x)_t)\lotimes_A M)\)
\par\noindent The final restriction formula on260 has an opening RΓ parenthesis with no terminal closing parenthesis after M. Miller silently closes it; retain the missing mark diplomatically and keep the closing mark in corrected forms.
\smallskip
\item[\hypertarget{apparatus:EN-0188}{188}] Page \hyperlink{printpage:en:264}{264}\quad\texttt{S45-SOL-1096}
\par\noindent\textbf{Printed:} \(\beta\circ \alpha = \beta\circ T^i(\alpha)\)
\par\noindent\textbf{Corrected:} \(\beta\circ \alpha = T^i(\beta)\circ\alpha\)
\par\noindent The scan visibly prints β∘T\textasciicircum{}i(α). With Hom\textasciicircum{}i(X,Y)=Hom(X,T\textasciicircum{}iY), the composable ordinary arrows are α followed by T\textasciicircum{}i(β), so correct to T\textasciicircum{}i(β)∘α.
\smallskip
\item[\hypertarget{apparatus:EN-0189}{189}] Page \hyperlink{printpage:en:265}{265}\quad\texttt{S45-ASTRA-V-0019}
\par\noindent\textbf{Printed:} \(F\circ T = T'\circ F\)
\par\noindent\textbf{Corrected:} \(F\circ T \simeq T'\circ F\)
\par\noindent The1976 footnote explicitly says that the original equality requirement was changed to an isomorphism. The unchanged equality sign in the displayed specification should be an isomorphism sign; the Codex proposal is accepted.
\smallskip
\item[\hypertarget{apparatus:EN-0190}{190}] Page \hyperlink{printpage:en:265}{265}\quad\texttt{S45-SOL-1100}
\par\noindent\textbf{Printed:} $u:X\leftrightarrow Y$, $v:Y\to Z$
\par\noindent\textbf{Corrected:} $u:X\to Y$, $v:Y\to Z$
\par\noindent The arrow in the printed list for u has arrowheads at both ends. The same page’s triangle and all ensuing axioms require the single morphism u:X→Y; retain the two-headed misprint diplomatically.
\smallskip
\item[\hypertarget{apparatus:EN-0191}{191}] Page \hyperlink{printpage:en:267}{267}\quad\texttt{S45-SOL-1108}
\par\noindent\textbf{Printed:} {\footnotesize\ttfamily\raggedright \textbackslash{}\allowbreak{}hom(\allowbreak{}X,\allowbreak{}M)\allowbreak{} \textbackslash{}\allowbreak{}ar[\allowbreak{}r]\allowbreak{}\textasciicircum{}\allowbreak{}-\allowbreak{}\{\allowbreak{}\textbackslash{}\allowbreak{}hom(\allowbreak{}w,\allowbreak{}M)\allowbreak{}\}\allowbreak{}}
\par\noindent\textbf{Corrected:} {\footnotesize\ttfamily\raggedright \textbackslash{}\allowbreak{}hom\textasciicircum{}\allowbreak{}i(\allowbreak{}X,\allowbreak{}M)\allowbreak{} \textbackslash{}\allowbreak{}ar[\allowbreak{}r]\allowbreak{}\textasciicircum{}\allowbreak{}-\allowbreak{}\{\allowbreak{}\textbackslash{}\allowbreak{}hom(\allowbreak{}w,\allowbreak{}M)\allowbreak{}\}\allowbreak{}}
\par\noindent The third term in the second exact sequence has Hom(X,M) without its degree superscript. The neighboring terms and degree-zero arrows show that the omitted superscript is i.
\smallskip
\item[\hypertarget{apparatus:EN-0192}{192}] Page \hyperlink{printpage:en:268}{268}\quad\texttt{S45-SOL-1112}
\par\noindent\textbf{Printed:} The morphisms of $\cA$ are
\par\noindent\textbf{Corrected:} The morphisms of $\eC(\cA)$ are
\par\noindent The second bullet on268 prints A alone. Morphisms of complexes are the morphisms of C(A), as both the immediately preceding definition and all subsequent constructions require.
\smallskip
\item[\hypertarget{apparatus:EN-0193}{193}] Page \hyperlink{printpage:en:269}{269}\quad\texttt{S45-SOL-1113}
\par\noindent\textbf{Printed:} additive automorphism of $\cA$
\par\noindent\textbf{Corrected:} additive automorphism of $\eC(\cA)$
\par\noindent The first prose paragraph on269 again has A alone. T is the translation automorphism on C(A), not on the original additive category A.
\smallskip
\item[\hypertarget{apparatus:EN-0194}{194}] Page \hyperlink{printpage:en:271}{271}\quad\texttt{S45-SOL-1120}
\par\noindent\textbf{Printed:} generated
    by its families
\par\noindent\textbf{Corrected:} generated
    by these families
\par\noindent The first bullet has ses familles. It points back to the three just-listed families of bounded complexes; the demonstrative ces is the intended reading.
\smallskip
\item[\hypertarget{apparatus:EN-0195}{195}] Page \hyperlink{printpage:en:271}{271}\quad\texttt{S45-SOL-1121}
\par\noindent\textbf{Printed:} We shall also pay for
\par\noindent\textbf{Corrected:} We shall also speak of
\par\noindent The scan unambiguously reads On paiera aussi. The paragraph introduces more subcategories rather than payment; parlera is the intended verb.
\smallskip
\item[\hypertarget{apparatus:EN-0196}{196}] Page \hyperlink{printpage:en:273}{273}\quad\texttt{S45-SOL-1124}
\par\noindent\textbf{Printed:} {\footnotesize\ttfamily\raggedright 0 \& d\_\allowbreak{}\{\allowbreak{}Y\textasciicircum{}\allowbreak{}\textbackslash{}\allowbreak{}bullet\}\allowbreak{}}
\par\noindent\textbf{Corrected:} {\footnotesize\ttfamily\raggedright 0 \& d\_\allowbreak{}\{\allowbreak{}Z\textasciicircum{}\allowbreak{}\textbackslash{}\allowbreak{}bullet\}\allowbreak{}}
\par\noindent The lower-right matrix entry visibly has Y. The decomposition Y=X+Z requires d\_Z in that corner. Preserve the printed Y as a mathematical misprint in diplomatic forms.
\smallskip
\item[\hypertarget{apparatus:EN-0197}{197}] Page \hyperlink{printpage:en:274}{274}\quad\texttt{S45-ASTRA-V-0021}
\par\noindent\textbf{Printed:} \(F^{i+-}(X)\)
\par\noindent\textbf{Corrected:} \(F^{i+1}(X)\)
\par\noindent A high-resolution scan crop shows a horizontal stroke after the plus, not a1. The exact sequence shifts cohomological degree by one. Accept the Codex correction; preserve the defective printed exponent diplomatically.
\smallskip
\item[\hypertarget{apparatus:EN-0198}{198}] Page \hyperlink{printpage:en:275}{275}\quad\texttt{S45-SOL-1126}
\par\noindent\textbf{Printed:} \(P^\bullet(Z^\bullet,f)\)
\par\noindent\textbf{Corrected:} \(P^\ast(Z^\bullet,f)\)
\par\noindent The parenthetical functor on complexes visibly has a bullet, although P\textasciicircum{}bullet is not defined until the ensuing quotient construction. At this point the functor is P\textasciicircum{}ast; the English already has the corrected symbol.
\smallskip
\item[\hypertarget{apparatus:EN-0199}{199}] Page \hyperlink{printpage:en:275}{275}\quad\texttt{S45-ASTRA-V-0039}
\par\noindent\textbf{Printed:} \(\hom^\circ:\eK(\cA)^\circ\)
\par\noindent\textbf{Corrected:} \(\hom^\bullet:\eK(\cA)^\circ\)
\par\noindent The enlarged275 crop shows the same hollow circle as the adjacent opposite-category superscript. Page309 explicitly refers back to this construction and writes Hom with a filled bullet, consistently with the complex-valued functor and its arguments. The distinct printed glyph is retained diplomatically rather than silently normalized.
\smallskip
\item[\hypertarget{apparatus:EN-0200}{200}] Page \hyperlink{printpage:en:281}{281}\quad\texttt{S45-SOL-1145}
\par\noindent\textbf{Printed:} let an object
    of $S_Y$
\par\noindent\textbf{Corrected:} let $t$ be an object
    of $S_Y$
\par\noindent The scan has et un objet de S\_Y, without naming that object. The subsequent source(t) and diagram show that the omitted variable is t; insert it only in corrected forms.
\smallskip
\item[\hypertarget{apparatus:EN-0201}{201}] Page \hyperlink{printpage:en:286}{286}\quad\texttt{S45-ASTRA-V-0026}
\par\noindent\textbf{Printed:} in which $f=i\circ h$ and $g\circ i=0$.
\par\noindent\textbf{Corrected:} in which $N'$ is an object of $\cB$, $f=i\circ h$, and $g\circ i=0$.
\par\noindent The scan omits N′∈B. Without it the condition is vacuous: choose N′=X, h=id\_X and i=f. The equivalence with surjectivity therefore requires this omitted membership hypothesis. Accept the later corrected-tree proposal.
\smallskip
\item[\hypertarget{apparatus:EN-0202}{202}] Page \hyperlink{printpage:en:288}{288}\quad\texttt{S45-ASTRA-V-0027}
\par\noindent\textbf{Printed:} The category $\cB^\bot/X$
\par\noindent\textbf{Corrected:} The category $X\backslash\cB^\bot$
\par\noindent The printed slash category repeats the notation for objects over X, whereas the same clause explicitly says under X and asks for an initial object. The corrected-tree replacement X\textbackslash{}B⊥ gives the comma category described by the words and the adjunction. Preserve the printed slash in diplomatic forms.
\smallskip
\item[\hypertarget{apparatus:EN-0203}{203}] Page \hyperlink{printpage:en:290}{290}\quad\texttt{S45-SOL-1169}
\par\noindent\textbf{Printed:} \(Q^\ast\circ Q\iso I^\bot\circ ^\ast i^\bot\)
\par\noindent\textbf{Corrected:} \(Q^\ast\circ Q\iso i^\bot\circ ^\ast i^\bot\)
\par\noindent The first textual isomorphism on290 visibly uses capital I\textasciicircum{}⊥. The inclusion functor was defined as lowercase i\textasciicircum{}⊥ on289, and the later triangle uses that lowercase symbol, so this is a print misprint.
\smallskip
\item[\hypertarget{apparatus:EN-0204}{204}] Page \hyperlink{printpage:en:290}{290}\quad\texttt{S45-SOL-1170}
\par\noindent\textbf{Printed:} {\footnotesize\ttfamily\raggedright \& \& I\textasciicircum{}\allowbreak{}\textbackslash{}\allowbreak{}bot\textbackslash{}\allowbreak{}circ \textasciicircum{}\allowbreak{}\textbackslash{}\allowbreak{}ast i\textasciicircum{}\allowbreak{}\textbackslash{}\allowbreak{}bot}
\par\noindent\textbf{Corrected:} {\footnotesize\ttfamily\raggedright \& \& i\textasciicircum{}\allowbreak{}\textbackslash{}\allowbreak{}bot\textbackslash{}\allowbreak{}circ \textasciicircum{}\allowbreak{}\textbackslash{}\allowbreak{}ast i\textasciicircum{}\allowbreak{}\textbackslash{}\allowbreak{}bot}
\par\noindent The lower-right vertex of the first diagram repeats the same capital I\textasciicircum{}⊥ misprint. Correct it consistently with the functor i\textasciicircum{}⊥.
\smallskip
\item[\hypertarget{apparatus:EN-0205}{205}] Page \hyperlink{printpage:en:290}{290}\quad\texttt{S45-SOL-1172}
\par\noindent\textbf{Printed:} \(\delta(TX_0)=-T\delta(X)\)
\par\noindent\textbf{Corrected:} \(\delta(TX)=-T\delta(X)\)
\par\noindent The final identity visibly adds a small round subscript to the left X. There is no distinguished object X\_0, and the right side is evaluated at arbitrary X. Remove the subscript in corrected forms while retaining the printed final δ notation.
\smallskip
\item[\hypertarget{apparatus:EN-0206}{206}] Page \hyperlink{printpage:en:292}{292}\quad\texttt{S45-SOL-1177}
\par\noindent\textbf{Printed:} \(\eK^{b,b}\)
\par\noindent\textbf{Corrected:} \(\eK^{b,b}(\cA)\)
\par\noindent The first natural-functor formula on292 omits(A) after K\textasciicircum{}\{b,b\}; the second formula and the definition on291 include it. Restore the argument only in corrected forms.
\smallskip
\item[\hypertarget{apparatus:EN-0207}{207}] Page \hyperlink{printpage:en:293}{293}\quad\texttt{S45-SOL-1180}
\par\noindent\textbf{Printed:} $m\geqslant 0$
\par\noindent\textbf{Corrected:} $m>0$
\par\noindent The print includes m=0 in the vanishing claim. At m=0 the identity of a nonzero object survives under the fully faithful embedding stated immediately above, so the correct range is m>0.
\smallskip
\item[\hypertarget{apparatus:EN-0208}{208}] Page \hyperlink{printpage:en:295}{295}\quad\texttt{S45-SOL-1186}
\par\noindent\textbf{Printed:} a morphism $f:X^\bullet \to Y^\bullet$ and
an isomorphism
\par\noindent\textbf{Corrected:} a morphism $f:X^\bullet \to Y^\bullet$ is
an isomorphism
\par\noindent The scan prints et un isomorphisme. The criterion is the assertion that f is an isomorphism, so est is required; preserve the conjunction in diplomatic French and English.
\smallskip
\item[\hypertarget{apparatus:EN-0209}{209}] Page \hyperlink{printpage:en:298}{298}\quad\texttt{S45-SOL-1194}
\par\noindent\textbf{Printed:} \(\ext^0(X^\bullet,Y^\bullet) \iso \varinjlim_{Y^\bullet/\mathsf{Qis}}\)
\par\noindent\textbf{Corrected:} \(\ext^0(X^\bullet,Y^\bullet) \iso \varinjlim_{Y^\bullet/\mathsf{Qis}^+}\)
\par\noindent In the bounded-below assertion the printed limit is indexed by all Qis. The companion bounded-above formula uses Qis\textasciicircum{}-; the corresponding strengthened bounded-below formula requires Qis\textasciicircum{}+. Retain the omitted plus diplomatically and restore it in corrected forms.
\smallskip
\item[\hypertarget{apparatus:EN-0210}{210}] Page \hyperlink{printpage:en:302}{302}\quad\texttt{S45-SOL-1207}
\par\noindent\textbf{Printed:} \(\h{}\circ \eK f\)
\par\noindent\textbf{Corrected:} \(\h^0{}\circ \eK f\)
\par\noindent The unbounded cohomological definition prints H∘Kf with no exponent. It is the zero-cohomology functor, as in the immediately preceding bounded definition; restore the omitted0 in corrected forms.
\smallskip
\item[\hypertarget{apparatus:EN-0211}{211}] Page \hyperlink{printpage:en:305}{305}\quad\texttt{S45-SOL-1221}
\par\noindent\textbf{Printed:} $Q\circ \Phi(Y)$
\par\noindent\textbf{Corrected:} $Q'\circ \Phi(Y)$
\par\noindent Item(c) prints Q∘Φ(Y), but Φ(Y) belongs to K(B) and must be sent to D(B) by Q′. Q has source K(A). The corrected-tree and English reading Q′ is justified by these explicitly given domains.
\smallskip
\item[\hypertarget{apparatus:EN-0212}{212}] Page \hyperlink{printpage:en:305}{305}\quad\texttt{S45-SOL-1222}
\par\noindent\textbf{Printed:} \(\varinjlim_{X/\mathsf{Qis}^\ast} Q\circ\Phi(-)\)
\par\noindent\textbf{Corrected:} \(\varinjlim_{X/\mathsf{Qis}^\ast} Q'\circ\Phi(-)\)
\par\noindent The right-derived limit formula likewise prints Q rather than Q′ before Φ. Its values lie in K(B), so the target-category quotient is Q′.
\smallskip
\item[\hypertarget{apparatus:EN-0213}{213}] Page \hyperlink{printpage:en:305}{305}\quad\texttt{S45-SOL-1223}
\par\noindent\textbf{Printed:} \(\varinjlim_{\mathsf{Qis}^\ast/X}\)
\par\noindent\textbf{Corrected:} \(\varprojlim_{\mathsf{Qis}^\ast/X}\)
\par\noindent The enlarged scan crop shows a right-pointing arrow beneath Lim in the left-derived formula. The dual construction is a limit over quasi-isomorphisms into X, hence an inverse limit; reverse the limit direction in corrected forms.
\smallskip
\item[\hypertarget{apparatus:EN-0214}{214}] Page \hyperlink{printpage:en:305}{305}\quad\texttt{S45-SOL-1224}
\par\noindent\textbf{Printed:} \(Q\circ \Phi(-)\)
\par\noindent\textbf{Corrected:} \(Q'\circ \Phi(-)\)
\par\noindent The left-derived formula also prints Q∘Φ. As in(c) and the right-derived formula, Q′ is required by the source and target categories.
\smallskip
\item[\hypertarget{apparatus:EN-0215}{215}] Page \hyperlink{printpage:en:309}{309}\quad\texttt{S45-ASTRA-V-0040}
\par\noindent\textbf{Printed:} \(\eK(\cA)\)
\par\noindent\textbf{Corrected:} \(\eK(\cA)^\circ\)
\par\noindent RHom(-,Y) is contravariant in X; the ensuing bifunctor has D(A)\textasciicircum{}circ explicitly. The scan omits the opposite-category superscript in this display. Accept the corrected-tree repair and propagate it to English.
\smallskip
\item[\hypertarget{apparatus:EN-0216}{216}] Page \hyperlink{printpage:en:310}{310}\quad\texttt{S45-SOL-1235}
\par\noindent\textbf{Printed:} {\footnotesize\ttfamily\raggedright 0 \textbackslash{}\allowbreak{}ar@\{\allowbreak{}<-\allowbreak{}\}\allowbreak{}[\allowbreak{}r]\allowbreak{}}
\par\noindent\textbf{Corrected:} {\footnotesize\ttfamily\raggedright 0 \textbackslash{}\allowbreak{}ar[\allowbreak{}r]\allowbreak{}}
\par\noindent The first arrow in the printed short exact sequence points left, from Y′ to0. Exactness in the stated order and the tensor sequence immediately below require0→Y′. Preserve the arrow reversal diplomatically.
\smallskip
\item[\hypertarget{apparatus:EN-0217}{217}] Page \hyperlink{printpage:en:311}{311}\quad\texttt{S45-SOL-1238}
\par\noindent\textbf{Printed:} {\footnotesize\ttfamily\raggedright Let \$X\$ be an object of \$\textbackslash{}\allowbreak{}eD\textasciicircum{}\allowbreak{}-\allowbreak{}(\allowbreak{}\textbackslash{}\allowbreak{}underline E)\allowbreak{}\$ and let \$Y\$ be an object of \$\textbackslash{}\allowbreak{}eD\textasciicircum{}\allowbreak{}+\allowbreak{}(\allowbreak{}\textbackslash{}\allowbreak{}underline\{\allowbreak{}E\}\allowbreak{}')\allowbreak{}\$. There is a functorial isomorphism \textbackslash{}\allowbreak{}[\allowbreak{}\textbackslash{}\allowbreak{}xymatrix\{\allowbreak{}   \textbackslash{}\allowbreak{}hom\_\allowbreak{}\{\allowbreak{}\textbackslash{}\allowbreak{}eD(\allowbreak{}\textbackslash{}\allowbreak{}underline\{\allowbreak{}E\}\allowbreak{}')\allowbreak{}\}\allowbreak{}(\allowbreak{}\textbackslash{}\allowbreak{}eL\textasciicircum{}\allowbreak{}-\allowbreak{} f\textasciicircum{}\allowbreak{}\textbackslash{}\allowbreak{}ast(\allowbreak{}X)\allowbreak{},\allowbreak{}Y)\allowbreak{} \textbackslash{}\allowbreak{}ar[\allowbreak{}r]\allowbreak{}\textasciicircum{}\allowbreak{}-\allowbreak{}\textbackslash{}\allowbreak{}sim     \& \textbackslash{}\allowbreak{}hom\_\allowbreak{}\{\allowbreak{}\textbackslash{}\allowbreak{}eD(\allowbreak{}\textbackslash{}\allowbreak{}underline E)\allowbreak{}\}\allowbreak{}(\allowbreak{}X,\allowbreak{}\textbackslash{}\allowbreak{}eR\textasciicircum{}\allowbreak{}+\allowbreak{} f\_\allowbreak{}\textbackslash{}\allowbreak{}ast(\allowbreak{}Y)\allowbreak{})\allowbreak{} \textbackslash{}\allowbreak{}text\{\allowbreak{}.\}\allowbreak{} \}\allowbreak{}\textbackslash{}\allowbreak{}]\allowbreak{}}
\par\noindent\textbf{Corrected:} {\footnotesize\ttfamily\raggedright Let \$X\$ be an object of \$\textbackslash{}\allowbreak{}eD\textasciicircum{}\allowbreak{}-\allowbreak{}(\allowbreak{}\textbackslash{}\allowbreak{}underline\{\allowbreak{}E\}\allowbreak{}')\allowbreak{}\$ and let \$Y\$ be an object of \$\textbackslash{}\allowbreak{}eD\textasciicircum{}\allowbreak{}+\allowbreak{}(\allowbreak{}\textbackslash{}\allowbreak{}underline E)\allowbreak{}\$. There is a functorial isomorphism \textbackslash{}\allowbreak{}[\allowbreak{}\textbackslash{}\allowbreak{}xymatrix\{\allowbreak{}   \textbackslash{}\allowbreak{}hom\_\allowbreak{}\{\allowbreak{}\textbackslash{}\allowbreak{}eD(\allowbreak{}\textbackslash{}\allowbreak{}underline E)\allowbreak{}\}\allowbreak{}(\allowbreak{}\textbackslash{}\allowbreak{}eL\textasciicircum{}\allowbreak{}-\allowbreak{} f\textasciicircum{}\allowbreak{}\textbackslash{}\allowbreak{}ast(\allowbreak{}X)\allowbreak{},\allowbreak{}Y)\allowbreak{} \textbackslash{}\allowbreak{}ar[\allowbreak{}r]\allowbreak{}\textasciicircum{}\allowbreak{}-\allowbreak{}\textbackslash{}\allowbreak{}sim     \& \textbackslash{}\allowbreak{}hom\_\allowbreak{}\{\allowbreak{}\textbackslash{}\allowbreak{}eD(\allowbreak{}\textbackslash{}\allowbreak{}underline\{\allowbreak{}E\}\allowbreak{}')\allowbreak{}\}\allowbreak{}(\allowbreak{}X,\allowbreak{}\textbackslash{}\allowbreak{}eR\textasciicircum{}\allowbreak{}+\allowbreak{} f\_\allowbreak{}\textbackslash{}\allowbreak{}ast(\allowbreak{}Y)\allowbreak{})\allowbreak{} \textbackslash{}\allowbreak{}text\{\allowbreak{}.\}\allowbreak{} \}\allowbreak{}\textbackslash{}\allowbreak{}]\allowbreak{}}
\par\noindent The scan reverses E and E′ in both object hypotheses and both Hom categories of the final adjunction. Since f*:E′→E and f\_*:E→E′, X belongs to D\textasciicircum{}-(E′), Y to D\textasciicircum{}+(E), the left Hom is in D(E), and the right Hom in D(E′). Correct all four positions consistently. The printed module-category notation is underlined E, restored in both branches.
\smallskip
\end{description}
\endgroup
\fi
